% concatenated from PFM_uniqueness_arxiv.tex
\UseRawInputEncoding % because problems with accents in refs



\documentclass[10pt]{article}
\usepackage[letterpaper, left=1in, top=1in, right=1in, bottom=1in, verbose, ignoremp]{geometry}
\usepackage[
appendix=inline    %%inline for proofs within paper, proofs in appendix by default
]{apxproof}

\usepackage{url}\RequirePackage[colorlinks,citecolor=blue, linkcolor=blue,urlcolor = blue]{hyperref}
\usepackage{latexsym,amssymb,amsmath,amsfonts,graphicx,color,fancyvrb,amsthm,enumerate,subcaption,mathrsfs}
\usepackage[longnamesfirst,authoryear,round]{natbib}
\usepackage{xy}\xyoption{all} \xyoption{poly} \xyoption{knot}
\usepackage{float}
\thispagestyle{empty} % no page number on front page
\usepackage{bm}
\usepackage{bbm}
\usepackage{multicol,multirow}
\usepackage[safe]{tipa}
\usepackage{array}
\usepackage{relsize}
\usepackage{chngcntr}
\usepackage{etoolbox}
\usepackage{caption}
\usepackage{tikz}
\usetikzlibrary{patterns}
\usepackage{tikz-3dplot}
\usetikzlibrary{decorations.pathreplacing,calc}
\usetikzlibrary{shapes,backgrounds}
\usetikzlibrary{patterns}
\usetikzlibrary{cd}
\usepackage{amsopn}
\usepackage{calc}
\usepackage{algorithm}
\usepackage[noend]{algpseudocode}
\usepackage{hyperref}
\usepackage{mathtools}
\newcommand{\indep}{\rotatebox[origin=c]{90}{$\models$}}
\usepackage{tikz-cd} % commutative diagrams
\usepackage{amscd} % commutative diagrams
\usepackage{comment}
\usepackage[capitalize,nameinlink,compress]{cleveref}
\crefname{figure}{Figure}{Figures} 
\crefname{equation}{}{} 
\crefname{assumption}{Assumption}{Assumptions}
\crefname{subsection}{Subsection}{Subsections}
\usepackage{textcomp}
\usepackage{mathabx}
\usepackage{thmtools}
\usepackage{thm-restate}

\newcounter{cdrow}
\newcommand\cdlabel[1]{%
	\Crefstepcounter{cdrow}
	\label{#1}
	(\thecdrow)
}

%%% Formatting theorems
\newtheorem{theorem}{Theorem}[section]
\newtheorem*{theorem*}{Theorem}
\newtheorem{corollary}[theorem]{Corollary}
\newtheorem{lemma}[theorem]{Lemma}
\newtheorem{proposition}[theorem]{Proposition}
\newtheorem{claim}[theorem]{Claim}
\newtheorem*{claim*}{Claim}

\theoremstyle{definition}
\newtheorem{definition}[theorem]{Definition}
\newtheorem*{definition*}{Definition}
\newtheorem{assumption}{Assumption}
\newtheorem{question}{Question}


\theoremstyle{remark}
\newtheorem{remark}[theorem]{Remark}
\newtheorem{remarks}[theorem]{Remarks}
\newtheorem{example}[theorem]{Example}
\newtheorem*{example*}{Example}

\DeclareMathOperator{\tr}{tr}


% custom definitions ...
\def\eq#1{(\Cref{#1})}
\def\pdf{p.d.f.\ } \def\cdf{c.d.f.\ }
\def\pdfs{p.d.f.s} \def\cdfs{c.d.f.s}
\def\mgf{m.g.f.\ } \def\mgfs{m.g.f.s\ }
\def\ci{\perp   \perp}  % conditional independence symbol
\def\diag{{\rm diag}}
\def\log{{\rm log}}
\def\tr{{\rm tr}}
\def\cond{\, | \,}
\newcommand\myeq{\stackrel{\mathclap{\normalfont\mbox{$t := r^p$}}}{=}}
\newcommand*\diff{\mathop{}\!\mathrm{d}}

\makeatletter
\newcommand*{\op}{%
	\DOTSB
	\mathop{\vphantom{\bigoplus}\mathpalette\matt@op\relax}%
	\slimits@
}
\newcommand\matt@op[2]{%
	\vcenter{\m@th\hbox{\resizebox{\widthof{$#1\bigoplus$}}{!}{$\boxplus$}}}%
}
\makeatother

\def\dsum{\displaystyle\sum}
\def\dint{\displaystyle\int}
\def\dsup{\displaystyle\sup}
\def\dinf{\displaystyle\inf}
\def\dmin{\displaystyle\min}
\def\dlim{\displaystyle\lim}
\def\bz{{\mathbf z}} % sample

\newcommand{\comp}{\mathbb{C}}
\newcommand{\sph}{\mathbb{S}}
\newcommand{\bola}{\mathbb{B}}
\newcommand{\disk}{\mathbb{D}}

\newcommand{\R}{\mathbb{R}}
\newcommand{\N}{\mathbb{N}}
\newcommand{\Z}{\mathbb{Z}}
\newcommand{\Ebb}{\mathbb{E}}
\newcommand{\Fbb}{\mathbb{F}}
\newcommand{\Kbb}{\mathbb{K}}
\newcommand{\Lbb}{\mathbb{L}}
\newcommand{\Mbb}{\mathbb{M}}
\newcommand{\Vbb}{\mathbb{V}}
\newcommand{\Xbb}{\mathbb{X}}
\newcommand{\Ybb}{\mathbb{Y}}
\newcommand{\Zbb}{\mathbb{Z}}

\newcommand{\prt}{\partial}

\newcommand{\Ker}{\mathrm{Ker} \,}
\renewcommand{\Im}{\mathrm{Im} \,}
\newcommand{\Coker}{\mathrm{Coker} \,}
\renewcommand{\H}{\mathrm{H}}
\newcommand{\rank}{\mathrm{rank} \,}
\renewcommand{\L}{\mathcal{L}}


\newcommand{\abf}{\mathbf{a}}
\newcommand{\cbf}{\mathbf{c}}
\newcommand{\fbf}{\mathbf{f}}

\newcommand{\VR}{\mathrm{VR}}

\newcommand{\ot}{\mathrm{OT}}
\newcommand{\T}{\intercal}
\newcommand{\Trop}{\mathrm{Trop}}
\newcommand{\Sym}{\mathrm{Sym}}
\newcommand{\argmin}{\text{argmin}}
\newcommand{\argmax}{\text{argmax}}
\newcommand{\codim}{\textrm{codim}}
\newcommand{\TPT}[1]{ \R^{#1}/\R \mathbf{1}}
\newcommand{\norm}[2]{\Vert #2 \Vert_{#1}}
\newcommand{\dtr}{d_{\mathrm{tr}}}
\newcommand{\esssupp}{\mathrm{esssupp}}
\newcommand{\CType}{F_1,\dots,F_n}
\newcommand{\sample}{ X_1,\dots, X_n}
\DeclareMathOperator{\tconv}{tconv}
\DeclareMathOperator{\conv}{conv}
\DeclareMathOperator{\relint}{relint}
\DeclareMathOperator{\closure}{Cl}
\DeclareMathOperator{\FMset}{\mathcal{F}}
\DeclareMathOperator{\aff}{aff}
\DeclareMathOperator{\proj}{proj}


\newcommand{\am}[1]{{\textcolor{cyan}{ #1  --Anthea}}} %%% Added by Anthea. For comments to appear in manuscript.
\newcommand{\rt}[1]{ {\textcolor{violet!60!magenta}{ #1  --Roan.}}} %%% Roan's comments
\newcommand{\amc}[1]{ {\textcolor{orange}{ #1  --Andy.}}} %%% Andy's comments


\newcommand{\ms}[1]{\mathscr{#1}}
\newcommand{\sgn}{\text{\sgn}}
\newcommand{\Vol}{\text{Vol}}
\newcommand{\one}{1}
\newcommand{\two}{2}
\newcommand{\id}{\text{id}}
\newcommand{\supp}{\text{Supp}}
\newcommand{\Poi}{\text{Poi}}
\newcommand{\Tan}{\text{Tan}}
\newcommand{\diam}{\text{diam}}
\newcommand{\Prob}{\mathbb{P}}
\newcommand{\E}{\mathbb{E}}
\newcommand{\W}{\mathcal{W}}
\renewcommand{\P}{\mathcal{P}}
\renewcommand{\vec}[1]{\mathbf{#1}}
\newcommand{\trop}{\text{trop}}
\newcommand{\im}{\text{Im}}
\renewcommand{\ker}{\text{Ker}}
\newcommand{\Cov}{\text{Cov}}
\newcommand{\Var}{\text{Var}}
\DeclareMathOperator{\type}{type}
\newcommand{\bl}[1]{{\mathbf #1}}
\newcommand{\bs}[1]{\boldsymbol #1}
\newcommand{\mb}[1]{\mathbb #1}
\newcommand{\mc}[1]{\mathcal #1}
\newcommand{\NR}{\mathcal{N}\mathcal{R}}
\newcommand{\A}{\mathcal{A}}

\makeatletter
\def\@biblabel#1{}
\makeatother

\makeatletter
\patchcmd{\NAT@citex}
{\@citea\NAT@hyper@{%
		\NAT@nmfmt{\NAT@nm}%
		\hyper@natlinkbreak{\NAT@aysep\NAT@spacechar}{\@citeb\@extra@b@citeb}%
		\NAT@date}}
{\@citea\NAT@nmfmt{\NAT@nm}%
	\NAT@aysep\NAT@spacechar\NAT@hyper@{\NAT@date}}{}{}

% Patch case where name and year are separated by opening bracket
\patchcmd{\NAT@citex}
{\@citea\NAT@hyper@{%
		\NAT@nmfmt{\NAT@nm}%
		\hyper@natlinkbreak{\NAT@spacechar\NAT@@open\if*#1*\else#1\NAT@spacechar\fi}%
		{\@citeb\@extra@b@citeb}%
		\NAT@date}}
{\@citea\NAT@nmfmt{\NAT@nm}%
	\NAT@spacechar\NAT@@open\if*#1*\else#1\NAT@spacechar\fi\NAT@hyper@{\NAT@date}}
{}{}

\makeatother
\usepackage[mathscr]{euscript}

%%%%%%%%%%%%%%%%%%%%%%%%%%%%%%%%%%%%%%%%%%%%%%%%%%%

\begin{document}
	
	\def\spacingset#1{\renewcommand{\baselinestretch}%
		{#1}\small\normalsize} \spacingset{1}
	%\spacingset{1.45} % DON'T change the spacing!
	
	
	\begin{flushleft}
		{\Large{\textbf{On the Uniqueness of Fr\'echet Means for Polytope Norms}}}
		\newline
		\\
		Roan Talbut$^{1,\dagger}$, Andrew McCormack$^{2}$, and Anthea Monod$^{1}$
		\\
		\bigskip
		\bf{1} Department of Mathematics, Imperial College London, UK
		\\
		\bf{2} Department of Mathematical and Statistical Sciences, University of Alberta, Edmonton, Canada
		\\
		\bigskip
		$\dagger$ Corresponding e-mail: roan.s.talbut@durham.ac.uk
	\end{flushleft}
	
	%%%%%%%%%%%%%%%%%%%%%%%%%%%%%%%%%%%%%%%%%%%%%%%%%%%
	
	\section*{Abstract}
	
	Fr\'echet means are a popular type of average for non-Euclidean datasets, defined as those points which minimise the average squared distance to a set of data points. We consider the behaviour of sample Fr\'echet means on normed spaces whose unit ball is a polytope; this setting is rarely covered by existing literature on Fr\'echet means, which focuses on smooth spaces or spaces with bounded curvature. We study the geometry of the set of Fr\'echet means over polytope normed spaces, with a focus on dimension and probabilistic conditions for uniqueness. In particular, we provide a geometric characterisation of the threshold sample size at which Fr\'echet means have a positive probability of being unique, and we prove that this threshold is at most one more than the dimension of our space. We are able to use this geometric characterisation to compute the unique Fr\'echet mean sample threshold in the case of the $\ell_\infty$ and $\ell_1$ norms.
	
	
	\paragraph{Keywords: polytope geometry, Fr\'echet means, polytope norms, uniqueness, small samples.}
	
	%%%%%%%%%%%%%%%%%%%%%%%%%%%%%%%%%%%%%%%%%%%%%%%%%%%
	
	\section{Introduction}
	The Fr\'echet mean set of a collection of data points in a metric space is defined as the set of points that minimise the average squared distance to the data points \citep{frechet1948elements}. Fr\'echet means generalise the classical Euclidean notion of a mean to datasets taking values in general metric spaces. As this definition only uses the metric of the underlying state space, Fr\'echet means have been applied to a wide variety of different geometric settings such as the space of PSD matrices \citep{arsigny2007geometric,fiori2009learning,petersen2019frechet}, the space of persistence diagrams in topological data analysis \citep{mileyko2011probability,cao2022approximating}, and Wasserstein space \citep{panaretos2020invitation,panaretos2020phase}. \cref{def:frechet_mean} gives the formal definition of Fr\'echet means on a general metric space, for both populations and samples.
	
	\begin{definition}\label{def:frechet_mean}
		Let $X$ be a random variable taking values in a general metric space $(\Omega,d)$. Its population Fr\'echet function $f: \Omega \rightarrow \R$ is defined as
		\[
		f(\vec \theta) = \E \left[ d(X,\vec \theta)^2 \right].
		\]
		For a sample $\{\vec x_1, \dots, \vec x_n\} \subset \Omega$, the sample Fr\'echet function is defined as
		\[
		f(\vec \theta) = \frac{1}{n}\sum_{i = 1}^n d(\vec x_i,\vec \theta)^2.
		\]
		The (population/sample) Fr\'echet mean set is then the set of global minima:
		\[
		\FMset = \argmin_{\vec \theta \in \Omega} f(\vec \theta).
		\]
		We often denote these Fr\'echet mean sets by $\FMset(X)$ or $\FMset(\vec x_1, \dots, \vec x_n)$.
	\end{definition}
	
	In this paper, we study the behaviour of sample Fr\'echet means on normed spaces whose unit ball is a polytope, hereafter referred to as spaces with a \emph{polytope norm}. This fundamental setting is not covered by most Fr\'echet mean literature, which instead focuses on manifolds \citep{BhatPatFMI,BhatPatFMII} or spaces with bounded curvature \citep{SturmHadamard,amendola2024invitation}. Despite this, polytope norms are prevalent in statistics; the $\ell_1$ norm is used in LASSO regression \citep{ranstam2018lasso} in order to encourage sparse parameterisation of regression models, and avoid overfitting to data. Our work directly addresses $\ell_1$ Fr\'echet means, as we show with running examples throughout this paper. However, our results are far more general, and apply to Fr\'echet means induced by any polytope norm. We also present examples using the $\ell_\infty$ norm to show the contrast in their behaviour.
	
	Fr\'echet means are defined as the solution set of an optimisation problem, so the question of existence and uniqueness is the natural first topic of study in any geometric setting. The existence of a Fr\'echet mean is equivalent to a distribution having finite moments, while the uniqueness of Fr\'echet means is more heavily dependent on the underlying geometry. \cite{SturmHadamard} showed that Fr\'echet means are unique in spaces with non-positive Alexandrov curvature, such as the cone of SPD matrices \citep{arsigny2007geometric}, and the BHV space of phylogenetic trees \citep{billera2001geometry,nye2017principal}. For Riemannian manifolds, \cite{BhatPatFMI} give sufficient conditions for uniqueness of sample and population Fr\'echet means, and more recently it was shown that for any absolutely continuous distribution over a complete Riemannian manifold, sample Fr\'echet means will almost surely be unique \citep{arnaudon2014meansUniqueness}. 
	
	In this paper, we present similarly foundational results on the uniqueness of Fr\'echet means from polytope norms. For data with an absolutely continuous underlying distribution, we show that the population Fr\'echet mean will be unique, while the probability of a unique sample Fr\'echet mean will converge to one. We then consider the probability of a unique Fr\'echet mean for small samples; we show that there is some sample threshold below which Fr\'echet means are almost never unique, and above which Fr\'echet means have a positive probability of uniqueness. This result follows from characterising the geometric configurations which can occur with positive probability; \cite{durier1985geometricalFermateber} take a similar view in their study of Fermat--Weber points, though their work is deterministic where ours is probabilistic. We show that this sample threshold is at most $k+1$ for general polytope norms, where $k$ is the dimension of our space. In the cases of the $\ell_\infty$ and $\ell_1$ norms, this threshold is $k$ and 3 respectively. 
	
	We provide a brief introduction to the relevant polytope geometry in \cref{sec:polytope_geometry} and the geometric foundations of Fr\'echet means on norm spaces in \cref{sec:Normed_Space_FMs}. We will show that for any absolutely continuous distribution, the population Fr\'echet mean will be unique, while the probability of a unique sample Fr\'echet mean is zero for small sample sizes and converging to one as the sample size increases (\cref{sec:large_sample_uniqueness}). Our main result, \cref{thm:wwp_dim_condition}, gives geometric conditions which dictate the number of samples required for this probability of uniqueness to go from zero to positive (\cref{sec:FM_uniqueness}). In \cref{sec:calculating_examples}, we present the threshold values for specific norms. We conclude by showing how the Frank--Wolfe algorithm \citep{frank1956algorithm} can be used for the exact computation of Fr\'echet means in \cref{sec:computation}.
	
	\paragraph{Notation.} We use $\relint(U), \closure(U)$ and $\aff(U)$ to denote the relative interior, closure and affine span of a set $U \subseteq \mb{R}^k$ respectively. We use $[k]$ to denote the set $\{1,\ldots,k\}$. We use $\vec e_i$ to denote the standard $i^{th}$ basis vecotr. Throughout this paper $f$ will refer to the sample Fr\'echet function unless otherwise specified.
	
	\section{Polytope Geometry} \label{sec:polytope_geometry}
	
	In this section, we review the definitions of polytopes, their faces and their polars. For a more detailed introduction, we refer the reader to \cite{ziegler2012lectures}. We then formalise a polytope norm --- a norm whose unit ball is a polytope --- and review the convex geometry corresponding to such a norm.
	
	\subsection{Polytopes}
	
	Polytopes can be viewed in two ways --- as a convex hull of finitely many points, or an intersection of finitely many halfspaces. 
	
	\begin{definition}[Polytope]
		A polytope $Q$ in $\R^k$ is the convex hull of a finite set of points $\{ \vec v_1, \dots, \vec v_r \} \subset \mathbb{R}^k$:
		\[
		Q = \conv(\{ \vec v_1, \dots, \vec v_r \}) = \left\{ \sum_{i=1}^r \lambda_i \vec v_i: \lambda_1, \dots, \lambda_r \geq 0, 
		\; \sum_{i=1}^r \lambda_i = 1 \right\}.
		\]
		Equivalently, given a matrix  $A \in \R^{r \times k}$ and a vector $\vec b \in \R^r$, a polytope can also be defined as the set of all solutions to the following linear inequalities:
		\[
		Q = P(A, \vec b) \coloneqq \left \{ \vec x \in \R^k \; |
		\; A \vec x \leq \vec b, \text{ coordinatewise} \right\},
		\]
		where it is assumed that the set of solutions is a bounded set. The convex hull and inequality representations of $Q$ are referred to as the vertex and halfspace representations respectively.
	\end{definition}
	
	Throughout this paper, we use two classes of polytopes as running examples; these are hypercubes and cross-polytopes, which arise as the unit balls of the $\ell_{\infty}$ and $\ell_1$ norms respectively. The vertex and halfspace representation of each of these balls is given in \Cref{ex:hypercube_representation,ex:cross-polytope_representation}, while \Cref{fig:polytope_pictures} shows visualisations of these polytopes for $k=2$ and 3.
	
	\begin{example}[Hypercubes]\label{ex:hypercube_representation}
		We define the $k$-dimensional hypercube $B_{\infty}^k$ by
		\[
		B_{\infty}^k \coloneqq \conv \left ( \{ -1, 1 \}^k\right) = \{ \vec x \in \R^k : \forall i \in [k], |x_i| \leq 1 \}.
		\]
	\end{example}
	\begin{example}[Cross-polytopes]\label{ex:cross-polytope_representation}
		We define the $k$-dimensional cross-polytope $B_1^k$ by 
		\[
		B_1^k \coloneqq \conv\left( \{ \pm \vec e_1, \dots, \pm \vec e_k \} \right) = \{ \vec x \in \R^k: \textstyle \sum_{i=1}^k |x_i| \leq 1\}.
		\]
	\end{example}
	
	\begin{figure*}[t!]
		\centering
		\begin{subfigure}[b]{0.4\textwidth}
			\centering
			\begin{tikzpicture}
				\filldraw[fill=gray,fill opacity=1/4] (1,1) -- (-1,1) -- (-1,-1) -- (1,-1) -- (1,1);
				\draw[->] (-2,0) -- (2,0);
				\draw[->] (0,-2) -- (0,2);
			\end{tikzpicture}
			\caption{2-dimensional hypercube}
		\end{subfigure}
		\begin{subfigure}[b]{0.4\textwidth}
			\centering
			\begin{tikzpicture}
				\filldraw[fill=gray,fill opacity=1/4] (1,0) -- (0,1) -- (-1,0) -- (0,-1) -- (1,0);
				\draw[->] (-2,0) -- (2,0);
				\draw[->] (0,-2) -- (0,2);
			\end{tikzpicture}
			\caption{2-dimensional cross-polytope}
		\end{subfigure}
		\begin{subfigure}[b]{0.4\textwidth}
			\centering
			\tdplotsetmaincoords{70}{50}
			\begin{tikzpicture}[scale=0.66, tdplot_main_coords]
				\coordinate (O) at (0,0,0);
				
				\draw[->] (-3,0,0) -- (3,0,0) node[anchor=north east]{$x$};
				\draw[->] (0,-3,0) -- (0,3,0) node[anchor=north west]{$y$};
				\draw[->] (0,0,-2.25) -- (0,0,2.25) node[anchor=south]{$z$};
				
				\coordinate (A) at (1,1,1);
				\coordinate (B) at (1,-1,1);
				\coordinate (C) at (-1,-1,1);
				\coordinate (D) at (-1,1,1);
				\coordinate (E) at (1,1,-1);
				\coordinate (F) at (1,-1,-1);
				\coordinate (G) at (-1,-1,-1);
				\coordinate (H) at (-1,1,-1);
				
				\draw (A) -- (B) -- (C) -- (D) -- (A);
				\draw (E) -- (F) -- (G);
				\draw[dashed] (G) -- (H);
				\draw[dashed] (E) -- (H);
				\draw (A) -- (E);
				\draw (B) -- (F);
				\draw (C) -- (G);
				\draw[dashed] (D) -- (H);
				\fill[fill=gray,opacity=1/8](A) -- (B) -- (F) -- (E) -- cycle;
				\fill[fill=gray,opacity=1/8](C) -- (D) -- (H) -- (G) -- cycle;
				\fill[fill=gray,opacity=1/8](A) -- (D) -- (H) -- (E) -- cycle;
				\fill[fill=gray,opacity=1/8](C) -- (B) -- (F) -- (G) -- cycle;
				\fill[fill=gray,opacity=1/8](A) -- (B) -- (C) -- (D) -- cycle;
				\fill[fill=gray,opacity=1/8](E) -- (F) -- (G) -- (H) -- cycle;
			\end{tikzpicture}
			\caption{3-dimensional hypercube}
		\end{subfigure}
		\begin{subfigure}[b]{0.4\textwidth}
			\centering
			\tdplotsetmaincoords{70}{50}
			\begin{tikzpicture}[tdplot_main_coords]
				\coordinate (O) at (0,0,0);
				
				\draw[->] (-2,0,0) -- (2,0,0) node[anchor=north east]{$x$};
				\draw[->] (0,-2,0) -- (0,2,0) node[anchor=north west]{$y$};
				\draw[->] (0,0,-1.5) -- (0,0,1.5) node[anchor=south]{$z$};
				
				\coordinate (A) at (1,0,0);
				\coordinate (B) at (0,1,0);
				\coordinate (C) at (0,0,1);
				\coordinate (D) at (-1,0,0);
				\coordinate (E) at (0,-1,0);
				\coordinate (F) at (0,0,-1);
				
				\draw (D) -- (E) -- (A) -- (B);
				\draw[dashed] (D) -- (B);
				\draw (D) -- (C) -- (A) -- (F);
				\draw[dashed] (D) -- (F);
				\draw (F) -- (E) -- (C) -- (B);
				\draw[dashed] (F) -- (B);
				\fill[fill=gray,opacity=1/8](A) -- (B) -- (C) -- cycle;
				\fill[fill=gray,opacity=1/8](D) -- (E) -- (F) -- cycle;
				\fill[fill=gray,opacity=1/8](A) -- (B) -- (F) -- cycle;
				\fill[fill=gray,opacity=1/8](C) -- (D) -- (E) -- cycle;
				\fill[fill=gray,opacity=1/8](A) -- (E) -- (C) -- cycle;
				\fill[fill=gray,opacity=1/8](D) -- (B) -- (F) -- cycle;
				\fill[fill=gray,opacity=1/8](A) -- (E) -- (F) -- cycle;
				\fill[fill=gray,opacity=1/8](D) -- (B) -- (C) -- cycle;
			\end{tikzpicture}
			\caption{3-dimensional cross-polytope}
		\end{subfigure}		
		\caption{The hypercubes and cross-polytopes in dimension 2 and 3.}
		\label{fig:polytope_pictures}
	\end{figure*}
	
	An important attribute of the geometry of a polytope is its \emph{facial} structure.  
	
	\begin{definition}[Face, face lattice]
		Let $\vec c \in \R^k$, $c_0 \in \R$ be such that $\vec c ^{\T} \vec x \leq c_0$ for all $\vec x$ in the polytope $Q$. A face $F$ of $Q$ is an intersection of the form:
		\[
		F= Q \bigcap \{ \vec x \in \R^k : \, \vec c^{\T} \vec x = c_0 \}.
		\]
		The set of all faces is denoted $\mathfrak{F}(Q)$, and forms a lattice with the partial ordering given by inclusion.
	\end{definition}
	
	We note that a polytope $Q$ is a face of itself, as is the empty set; any other face of $Q$ is referred to as a \emph{proper} face of $Q$. The $0$-dimensional faces of a polytope are called vertices, the $1$-dimensional faces of a polytope are edges, and the maximal proper faces of $Q$ are \emph{facets}. When $Q$ is given by $\conv \{ \vec v_1, \dots, \vec v_r \}$, every face of $Q$ will be given by the convex hull of some subset of $\{ \vec v_1, \dots, \vec v_r \}$. When $Q$ is given by solutions to a collection of linear inequalities, every face of $Q$ is the solution set of the same inequalities, with some taken to be equalities.
	
	Lastly, we introduce the polar of a polytope which is intimately tied to subgradients of polytope norms.
	
	\begin{definition}[Polar Polytope]
		Let $Q = \conv \{ \vec v_1, \dots, \vec v_r \} \subset \R^k$ be a full-dimensional polytope containing the origin in its interior. The polar of $Q$ is the polytope given by:
		\[
		Q^{\Delta} \coloneqq \{ \vec x \in \R^k : \forall \vec y \in Q , \, \vec y^{\T} \vec x \leq 1 \} = P(V,\mathbf 1),
		\]
		where $V$ is the matrix with rows $\vec v_1, \dots, \vec v_r$. The faces of $Q^{\Delta}$ are given by $\{ F^{\Diamond}: \, F \in \mathfrak{F}(Q) \}$, where
		\[
		F^{\Diamond} \coloneqq \{ \vec x \in Q^{\Delta} : \forall \vec y \in F, \; \vec y^{\T} \vec x = 1 \}.
		\]
		We call $F^{\Diamond}$ the polar of the face $F$.
	\end{definition}
	
	The faces of a polytope are in bijection with the faces of its polar polytope, in fact their face lattices are anti-isomorphic; $F_1^{\Diamond} \subseteq F_2^{\Diamond}$ if and only if $F_2 \subseteq F_1$. The polar of a polar polytope $Q^{\Delta \Delta}$ is equal to $Q$, and as such, we see that the polar of $P(A, \vec 1)$ is given by the convex hull of the row vectors of $A$. 
	
	\begin{example}[Hypercubes and Cross-polytopes]
		The $k$-dimensional hypercube and cross-polytope are polars of each other. The polars of the hypercube vertices are the facets of the cross-polytope, and vice versa.
	\end{example}
	
	\subsection{Polytope Norms}
	
	We now review the convex analysis of a norm whose unit ball is a polytope $B$. When defining polytope norms, the hyperplane representation is most natural.
	
	\begin{definition}[Polytope Normed Space]
		Let $B = P(A, \vec 1) \subset \R^k$ be a full-dimensional, centrally symmetric polytope, meaning that $\vec x \in B$ if and only if $- \vec x \in B$. Then $(\R^k, \norm{B}{\cdot})$ is a polytope normed space where the $B$-norm is defined by
		\begin{align}
			\label{eqn:BnormDefinition}
			\norm{B}{\vec x} \coloneqq \inf \{ d \in \R_{>0}: \,  d^{-1}\vec x \in B \} \equiv \max_{i \in [r]} (A \vec x)_i.
		\end{align}
		The metric induced by the $B$-norm is denoted by $d_B$. 
	\end{definition}
	
	It is straightforward to show that $\norm{B}{\vec x}$ is a norm, and that the unit ball of $\norm{B}{\cdot}$ is the polytope $B$. In the case of our standard polytope examples:
	
	\begin{itemize}
		\item The hypercube $B_{\infty}^k$ gives rise to the $\ell_{\infty}$ norm $\norm{B_{\infty}^k}{\vec x} \equiv \|\vec x \|_{\infty} = \max_{i \in [k]} \left \{ |x_i| \right \}$.
		\item The cross-polytope $B_{1}^k$ gives rise to the $\ell_{1}$ norm $\norm{B_{1}^k}{\vec x} \equiv \|\vec x \|_{1} = \sum_{i \in [k]} |x_i|$.
	\end{itemize}
	
	To discuss the local behaviour of a polytope norm near $\vec x$, we need to know which rows of $A$ are active, meaning that they attain the componentwise maximum of $A\vec x$. We refer to these rows as the \emph{active constraints}. In the following definition, we write $\vec a \in A$ to mean $\vec a$ is a row vector of $A$.
	
	\begin{definition}
		Let $F$ be a proper face of $B$. The active constraints on $F$, denoted by $\A(F)$, is the set:
		\[
		\A(F) = \{ \vec a \in A : \text{ for all } \vec x \in F, \,  \vec a^{\T} \vec x=1 \}. 
		\]
		Similarly, for any $\vec x \in \R^k$ we define $\A(\vec x)$ to be the active constraints of $\| \vec x \|_B$. That is:
		\[
		\A( \vec x ) = \{ \vec a \in A :  \vec a^{\T} \vec x = \| \vec x \|_B \}.
		\]
	\end{definition}
	Notice that $\conv \A(F) = F^{\Diamond}$. We also note that the set of open cones generated by proper faces of $B$, $\relint (\R_+ F)$, partition $\R^k \setminus \{ \vec 0 \}$ according to the active constraints; $\vec x \in \relint (\R_+ F)$ if and only if $\A(\vec x) = \A(F)$.  This partition of the space is the \emph{face fan} of $B$.
	
	Finally, we review the subgradient and directional derivatives of the $B$-norm at $\vec x$. 
	
	\begin{definition}[Subgradient, Directional Derivative]
		The subgradient of a convex function $g:\R^k \rightarrow \R$ at $\vec x$ is given by:
		\[
		\partial g(\vec x) \coloneqq \{ \vec v : \, g(\vec z) \geq g(\vec x) + \vec v^{\T}(\vec z-\vec x) \; \text{ for all } \vec z \in \R^k \}.
		\]
		The directional derivative in direction $\vec t \in \R^k$, denoted by $\partial_{\vec t} g(\vec x)$, is 
		\[
		\partial_{\vec t} g(\vec x) \coloneqq \lim_{h \rightarrow 0^+} h^{-1} \left( g(\vec x + h \vec t) - g(\vec x) \right) = \sup_{\vec v \in \partial g(\vec x)} \vec v^{\T} \vec t.
		\]
		This equality is from Theorem 23.4 of \citep{rockafellar1970convex}.
	\end{definition}
	
	
	\begin{remark}
		\label{rmk:subgradient}
		The subgradient of a maximum of differentiable functions $g(\vec x) = \max \{ g_1(\vec x), \dots, g_n(\vec x)\}$ is known to be the convex hull of the gradients of active functions at $\vec x$: $\partial g(\vec x) = \conv \{ \nabla_{\vec x} g_i(\vec x):g_i(\vec x) = g(\vec x)\}$. Applying this result to a polytope $B$-norm, which is a maximum of linear functions by \eqref{eqn:BnormDefinition}, we obtain that
		\begin{align*}
			\partial \Vert \vec x \Vert_B^2 =  \conv \left \{ \nabla_{\vec x} (\vec a^{\top}\vec x)^2: \vec a \in F^{\Diamond} \right\} = 2\| \vec x \| F^{\Diamond}, 
		\end{align*}
		when $\vec x \in \relint (\R_+ F)$. Consequently, the subgradient of the sample Fr\'echet function will be given by
		$\partial f(\vec x) =   2\sum_{i \in [n]} \| \vec x_i - \theta \| F_i^{\Diamond}$, where $F_i$ is the unique face such that $\vec x_i - \theta \in \relint (\R_+ F_i)$, and the summation is a Minkowski sum. Since a convex function is minimized if and only if its subgradient contains $\vec 0$, it follows that $\theta$ is a sample Fr\'echet mean if and only if $\vec 0 \in  \sum_{i \in [n]} \| \vec x_i - \theta \| F_i^{\Diamond}$.  
	\end{remark}
	
	\section{Fr\'echet Means on Normed Spaces} \label{sec:Normed_Space_FMs}
	
	In this section, we discuss Fr\'echet means and outline their fundamental properties in the case of (polytope) normed spaces. All Fr\'echet means in the sections to follow are taken with respect to a given $B$-norm, unless stated otherwise. As defined in \cref{def:frechet_mean}, we use $f$ to denote a Fr\'echet function; this is given by $f(\theta) = \E [d(X,\theta)^2]$ for general random variables and $f(\theta)=n^{-1} \sum_{i=1}^n d(\vec x_i,\theta)^2$ for a sample $\{ \vec x_1, \dots, \vec x_n \}$. We then denote the Fr\'echet mean set by $\FMset \coloneqq \argmin_{\theta} f(\theta)$, which we may abbreviate to $\FMset (X)$ or $\FMset(\vec x_1, \dots, \vec x_n)$.
	
	Understanding the convexity of the Fr\'echet function is crucial for understanding the existence and convexity of Fr\'echet mean sets; these foundational convexity results have been established in many settings, such as Wasserstein spaces \citep{panaretos2020frechet}, Riemannian spaces \citep{afsari2011riemannian}, metric trees, \citep{romon2023convex}, and the tropical projective torus \citep{lin2025tropical}. Below, we state the convexity results we require in the case of a general normed space. These results below are generalisations of the convexity results established in the tropical setting by \cite{lin2025tropical}, and mirror the work of \citep{durier1985geometricalFermateber} on Fermat--Weber points.
	
	\begin{lemma}
		\label{lemma:NormConvexity}
		Fix a normed space $(\R^k, \norm{}{\cdot})$ with induced metric $d$. For all $\vec x$, $d(\cdot, \vec x)^2$ is convex, and in fact if, for some $t \in (0,1)$ we have $$d((1-t)\vec\theta_1 + t \vec\theta_2, \vec x)^2  = (1-t)d(\vec\theta_1, \vec x)^2+ t d(\vec\theta_2, \vec x)^2,$$ then $d(\vec\theta_1, \vec x) = d(\vec\theta_2, \vec x)$. The Fr\'echet function is also convex, and so Fr\'echet mean sets are convex.
	\end{lemma}
	\begin{proof}
		Convexity of a squared norm is immediate from the convexity of norms, as the function $y \mapsto y^2$ is strictly monotone increasing and strictly convex over the domain $\R_{\geq 0}$:
		\begin{align*}
			\Vert (1-t)\vec\theta_1 + t \vec\theta_2 - \vec x\Vert^2 &\leq \left( (1-t)\Vert \vec\theta_1 - \vec x \Vert + t \Vert \vec\theta_2 - \vec x\Vert \right)^2, \\
			&\leq (1-t)\Vert \vec \theta_1 - \vec x \Vert^2 + t \Vert \vec \theta_2 - \vec x\Vert^2.
		\end{align*}
		Hence $\Vert \theta - x\Vert^2 = d(\theta,\vec x)^2$ is strictly convex in $\theta$. \\
		The second statement follows from the strict convexity of $y \mapsto y^2$; we have equality if and only if $\Vert \theta_1 - \vec x \Vert^2 = \Vert \theta_2 - \vec x\Vert^2$. The convexity of the Fr\'echet function $f$ is a direct result of expectation preserving convexity, while the convexity of Fr\'echet mean sets follows from the convexity of $f$.
	\end{proof}
	
	Building on \Cref{lemma:NormConvexity}, we will see that Fr\'echet mean sets are given by intersections of spheres. \Cref{lemma:FM_dists} expresses this in the case of a general population $X \sim \mb P$, while \Cref{cor:FMisPolytope} considers the case of a finite sample.
	
	\begin{lemma}
		\label{lemma:FM_dists}
		Let $X \sim \mathbb P$ be a random variable in $(\R^k, \| \cdot \|)$. Assume that $\theta_1$ and $\theta_2$ are Fr\'echet means of $X$. Then for every $\vec x$ in the  $\esssupp(\mathbb P)$, we have $d(\vec x, \theta_1) = d(\vec x, \theta_2)$. 
	\end{lemma}
	\begin{proof} 
		Fix some $\vec x \in \esssupp(\mb P)$. By definition, for all $m \in \N$ we have $\mb P(B_{1/m}(\vec x)) > 0$. We will show that there is some $\vec x_m \in B_{1/m}(\vec x)$ such that $d(\vec x_m,\theta_1) = d(\vec x_m, \theta_2)$ by contradiction. Let $\theta_3 = \tfrac{1}{2}(\theta_1 + \theta_2)$ be a midpoint of $\theta_1, \theta_2$. By the convexity of $d(X,\cdot)$ and strict convexity of $y \mapsto y^2$:
		\begin{align*}
			d(X,\theta_3)^2 \leq \big(\tfrac{1}{2}d(X,\theta_1) + \tfrac{1}{2}d(X,\theta_2)\big)^2 \leq \tfrac{1}{2}\big( d(X,\theta_1)^2 + d(X,\theta_2)^2\big).
		\end{align*}
		We assume for all $\vec u \in B_{1/m}(\vec x)$, $d(\vec u,\theta_1) \neq d(\vec u, \theta_2)$, so the above inequality is strict. It follows that
		\begin{align*}
			\E \left[I_{X \in B_{1/m}}d(X,\theta_3)^2\right] &< \tfrac{1}{2}\left(\E\left[I_{X \in B_{1/m}}d(X,\theta_1)^2\right] + \E \left[I_{X \in B_{1/m}}d(X,\theta_2)^2\right]\right).
		\end{align*}
		Hence
		\[
		\E \left[d(X,\theta_3)^2\right] < \tfrac{1}{2}\left(\E\left[d(X,\theta_1)^2\right] + \E \left[d(X,\theta_2)^2\right]\right),
		\]
		which contradicts the optimality of $\theta_1$ and $\theta_2$. Therefore for all $m \in \N$, we can pick some $\vec x_m \in B_{1/m}(\vec x)$ such that $d(\vec x_m,\theta_1) = d(\vec x_m,\theta_2)$. By continuity of the distance function we conclude $d(\vec x,\theta_1) = d(\vec x,\theta_2)$.
	\end{proof}
	
	As a natural corollary to \Cref{lemma:FM_dists}, we see that any distribution whose essential support is $\R^k$ must have a unique Fr\'echet mean (\cref{thm:unique_population_FM}). However, for finite samples, uniqueness is not guaranteed. Fr\'echet means are given by intersections of spheres centred at each data point, as shown in \Cref{cor:FMisPolytope}. In fact, in the case of polytope norms, Fr\'echet means are intersections of faces. \Cref{fig:FM_sets} provides a visualisation of this behaviour for Fr\'echet mean sets of a small number of samples. 
	
	\begin{corollary}
		\label{cor:FMisPolytope}
		Fix a normed space $(\R^k, \| \cdot \|_B)$ with unit ball $B$. There exist unique $d_i \in \R_{\geq 0}$ such that $$\FMset (\vec x_1,\ldots,\vec x_n) = \bigcap_{i \in [n]}\{d_iB + \vec x_i\} = \bigcap_{i \in [n]}\{d_i\partial B + \vec x_i\},
		$$ where $\partial B$ is the boundary of $B$. In particular, if $B$ is a polytope then $\FMset(\vec x_1, \dots, \vec x_n)$ is a polytope. Moreover, for every $i \in [n]$ there exists a proper face $F_i$ of $B$ such that $\FMset$ is given by $\bigcap_{i \in [n]}\{d_iF_i + \vec x_i\}$. 
	\end{corollary}
	
	\begin{proof}
		By \Cref{lemma:FM_dists}, the distances $d_i = d(\vec x_i,\theta)$ are fixed for all $\theta \in \FMset(\vec x_1, \dots, \vec x_n)$. Hence:
		\[
		\FMset(\vec x_1, \dots, \vec x_n) \subseteq \bigcap_{i \in [n]} \{d_i\partial B + \vec x_i\} \subseteq \bigcap_{i \in [n]} \{d_i B + \vec x_i\} \subseteq \FMset(\vec x_1, \dots, \vec x_n).
		\]
		This follows from \Cref{lemma:FM_dists}, nested intersections, and the definition of the $d_i$. When $B$ is a polytope, the intersection of balls about each $\vec x_i$ is also a polytope. 
		
		We prove the last claim by contradiction. Suppose there is no constraint $\vec a \in A$ which is active for all $\{ \vec x_i - \theta: \theta \in \FMset\}$. Then for all $\vec a_j \in A$, we have some $\theta_1, \dots, \theta_r \in \FMset$ such that $\vec a_j \in A$ is not in $\A(\vec x_i - \theta_j)$. By the convexity of $\FMset$, $\tfrac{1}{r}\sum_{j \in [r]} \theta_j$ is a Fr\'echet mean and must be in the set $d_i\partial B + \vec x_i$. For all $\vec a_j$, we have
		\[
		\vec a_j^\T(\vec x_i - \theta_j) < \Vert \vec x_i - \theta_j\Vert_B = d_i, \qquad \vec a_j^\T(\vec x_i - \theta_k) \leq \Vert \vec x_i - \theta_k\Vert_B = d_i.
		\]
		We obtain that
		\begin{align*}
			\left \Vert \vec x_i - \tfrac{1}{r}\textstyle\sum_{j \in [r]} \theta_j \right \Vert_B = \textstyle\max_{j'} \vec a_{j'}^{\top} \left(\tfrac{1}{r}\sum_{j \in [r]} (\vec x_i -  \theta_j) \right) < \tfrac{1}{r}\sum_{j \in [r]} \|\vec x_i -  \theta_j\|_B = d_i,
		\end{align*}
		contradicting $\tfrac{1}{r}\sum_{j \in [r]} \theta_j \in \FMset$. So for all $\vec x_i$, there is some face $F_i$ of $B$ with $\FMset \subseteq d_i F_i + \vec x_i$. Hence:
		\[
		\FMset(\vec x_1, \dots, \vec x_n) \subseteq \bigcap_{i \in [n]} \{d_iF_i + \vec x_i\} \subseteq \bigcap_{i \in [n]} \{d_i\partial B + \vec x_i\} = \FMset(\vec x_1, \dots, \vec x_n).
		\]
		So the result is proved.
	\end{proof}
	
	We have proved a more general result in the final part of this corollary, which will be useful later:
	
	\begin{corollary}\label{cor:convex_disjoint_from_relint}
		Suppose $S$ is a convex set contained in a polytope $Q$. If $S$ and $\relint Q$ are disjoint, then $S$ is contained in some proper face of $Q$.
	\end{corollary}
	
	\begin{figure}[!ht]
		\centering
		\begin{subfigure}{0.49\textwidth}
			\centering
			\begin{tikzpicture}[scale=1.8]
				\draw (1,1) -- (-1, 1) -- (-1,-1) -- (1,-1) -- (1,1);
				\draw (1,0) -- (1,2) -- (3,2) -- (3,0) -- (1,0);
				\filldraw (0,0) circle (1pt) node[anchor=north] {\small$\vec x_1 = (0,0)$};
				\filldraw (2,1) circle (1pt)  node[anchor=south] {\small$\vec x_2 = (2,1)$};
				\draw[red, line width = 2pt]  (1,0) -- (1,1);
			\end{tikzpicture}
		\end{subfigure}
		\hfill
		\begin{subfigure}{0.49\textwidth}
			\centering
			\begin{tikzpicture}[scale=0.9]
				\draw (0,0) -- (-4, 0) -- (-4,-4) -- (0,-4) -- (0,0);
				\draw (-0.5,0) -- (1.5, 0) -- (1.5,2) -- (-0.5,2) -- (-0.5,0);
				\draw (0,0.5) -- (0,-1.5) -- (2,-1.5) -- (2,0.5) -- (0,0.5);
				\filldraw (-2,-2) circle (1pt) node[anchor=north east] {\small$\vec x_1 = (-4,-4)$};
				\filldraw (0.5,1) circle (1pt)  node[anchor=south] {\small$\vec x_2 = (1,2)$};
				\filldraw (1,-0.5) circle (1pt)  node[anchor=north west] {\small$\vec x_3 = (2,-1)$};
				\filldraw[red] (0,0) circle (3pt);
			\end{tikzpicture}  
		\end{subfigure}
		\caption{Fr\'echet mean sets (red) are contained in faces of balls that are centred about each data point (\Cref{cor:FMisPolytope}).}
		\label{fig:FM_sets}
	\end{figure}
	
	\begin{remark}
		Corollary \ref{cor:FMisPolytope} highlights why Fr\'echet means with respect to polytope balls are particularly interesting to study, as such balls often have intersections that are positive dimensional, as is illustrated in the left hand side of Figure \ref{fig:FM_sets}. By contrast, the Fr\'echet mean of two points $\vec x_1, \vec x_2$ with respect to a norm with a strictly convex ball will always be the unique point where the balls around $\vec x_1$ and $\vec x_2$ have a single contact point with each other.
	\end{remark}
	
	We conclude with a final corollary on Fr\'echet means for normed spaces with strictly convex balls, to show the contrast with the polytope norm setting; if the unit ball of a norm is strictly convex, then Fr\'echet means must be unique.
	
	\begin{corollary}
		If $B$ is a strictly convex ball, in the sense that the relative interior of any line between two points in $\partial B$ is contained in $\mathrm{int}(B)$, then $\FMset (\vec x_1,\ldots,\vec x_n)$ consists of a single point. In particular, sample Fr\'echet means with respect to $\ell_p$ norms, $p \in (1,\infty)$, are always unique.
	\end{corollary}
	\begin{proof}
		The norm of a strictly convex ball has the property that whenever $\vec x$ is not proportional to $\vec y$ with respect to a non-negative scaling, we have for all $t \in (0,1)$ that 
		\begin{align*}
			\Vert (1-t)\vec x + t \vec y\Vert_B &= \big((1-t)\Vert \vec x \Vert_B + t\Vert \vec y \Vert_B\big)     \bigg\Vert \frac{(1-t)\Vert \vec x \Vert_B}{(1-t)\Vert \vec x \Vert_B + t\Vert \vec y \Vert_B} \frac{\vec x}{\Vert \vec x \Vert_B} + \frac{t\Vert \vec y \Vert_B}{(1-t)\Vert \vec x \Vert_B + t\Vert \vec y \Vert_B} \frac{\vec y}{\Vert \vec y \Vert_B} \bigg\Vert_B 
			\\
			&< (1-t)\Vert \vec x\Vert_B + t \Vert \vec y\Vert_B,
		\end{align*}
		since the above expression is a convex combination of the points $\vec x/\Vert \vec x \Vert_B, \vec y/\Vert \vec y \Vert_B \in \partial B$. Consider two Fr\'echet means, $\theta_1,\theta_2 \in \FMset(\vec x_1, \dots, \vec x_n)$. By the convexity of $\FMset$, and \cref{lemma:FM_dists}, $(1-t)\theta_1 + t\theta_2$ has the same distance to $\vec x_1$ as $\theta_1$ and $\theta_2$. So by the result above we must have that $\vec x_1 - \theta_1 = c(\vec x_1 - \theta_2)$ for some $c \geq 0$, as well as $d(\vec x_1, \theta_1) = d(\vec x_1, \theta_2)$. The only way this can occur is if $c = 1$, and thus $\theta_1 = \theta_2$. 
		
		The $\ell_p$ balls are strictly convex for $p \in (1,\infty)$ since if $\Vert \vec x \Vert_p = \Vert \vec y \Vert_p = 1$ with $\vec x \neq \vec y$, then  
		\begin{align*}
			\Vert (1-t)\vec x + t \vec y \Vert_p^p =   \textstyle\sum_{i = 1}^k \vert (1-t)x_i + ty_i\vert^p & \leq     \textstyle\sum_{i = 1}^k  ((1-t)\vert x_i \vert + t\vert y_i \vert)^p
			\\
			& <   \textstyle\sum_{i = 1}^k  \big((1-t)\vert x_i \vert^p  + t\vert y_i \vert^p\big) = 1.
		\end{align*}
		The strict inequality follows from Jensen's inequality and the fact that $\vec x \neq \vec y$.
	\end{proof}
	
	\section{Large Sample Uniqueness} \label{sec:large_sample_uniqueness}
	
	We continue our study of polytope Fr\'echet means by determining the probability that a sample Fr\'echet mean is unique, assuming that distribution of the samples has a positive density on $\R^k$ with respect to the Lebesgue measure.
	
	\begin{theorem}\label{thm:unique_population_FM}
		Suppose $X \sim \mb P$, with $\esssupp (\mb P) = \R^k$. Then $\FMset (X)$ is unique.
	\end{theorem}
	
	\begin{proof}
		This follows directly from \Cref{lemma:FM_dists}. 
	\end{proof}
	
	As shown by \cite{BhatPatFMI}, sample Fr\'echet means are strongly consistent estimators for population Fr\'echet means.
	
	\begin{theorem}\label{thm:consistent_FM}
		Let $\{X_i\}_{i \in \N}$ be a sequence of i.i.d. random vectors with $X_i \sim \mb{P}$, where $\mb P$ has a unique population Fr\'echet mean $\theta_*$. Let  $\theta_n$ be a random variable such that $\theta_n \in \FMset(X_1, \dots, X_n)$ almost surely. Then $\theta_n$ converges almost surely to $\theta_*$.
	\end{theorem}
	
	\begin{proof}
		This is a direct application of Theorem 2.3 of \citep{BhatPatFMI}. We note that closed and bounded subsets of $\R^k$ are compact, and we have assumed the population Fr\'echet function is finite. We have also assumed that $\mb P$ has a unique Fr\'echet mean $\theta_*$. Then by Theorem 2.3 of \cite{BhatPatFMI}, $\theta_n$ is a strongly consistent estimator for $\theta_*$.
	\end{proof}
	
	The consistency of sample Fr\'echet means does not guarantee their uniqueness. In \cref{thm:uniqueness_prob_convergence}, we show that the probability of a unique sample Fr\'echet mean tends to one as the sample size increases.
	
	\begin{theorem}\label{thm:uniqueness_prob_convergence}
		Suppose $X \sim \mb P$, with $\esssupp(\mb P) = \R^k$. Let $p_n$ be the probability that $\FMset(X_1,\dots,X_n)$ consists of a single element. Then $p_n$ converges to 1 as $n \rightarrow \infty$.
	\end{theorem}
	\begin{proof}
		Suppose that there is some sample Fr\'echet mean $\theta \in \FMset(\sample)$ such that for every possible active constraint $\vec a_i \in A$, there is some data point $X_{j(i)}$ satisfying $\A(X_{j(i)}-\theta) = \vec a_i$. Then by \Cref{cor:FMisPolytope}, the Fr\'echet mean set is contained in the intersection of corresponding faces, $\bigcap_i \|\theta-X_{j(i)}\| F_i+X_{j(i)}$ and so
		\begin{align}
			\FMset(X_1,\dots,X_n) \subseteq \{ \vec x: \vec a_i^{\T}(\vec x-X_{j(i)}) = \|\theta-X_{j(i)}\| \} \label{eq:FM_uniqueness_linear_constraints}.
		\end{align}
		As the set of all $\vec a_i$ spans $\R^k$, the linear constraints in \ref{eq:FM_uniqueness_linear_constraints} define a unique point in $\FMset(X_1, \dots, X_n)$. It now suffices to show that, with probability tending to 1, there is some Fr\'echet mean $\theta$ such that $A = \{ \A(X_j - \theta), j=1,\dots,n \}$.
		
		For each facet $F_i$, consider the intersection $S_i = \bigcap_{\theta \in B_1(\theta_*)} (\theta - \relint (\R_+ F_i)) \subset \R^k$ (see \Cref{fig:unique_prob_to_1}). This is the set of points $\vec x$ with $\A(\vec x - \theta) = \vec a_i$ for all $\theta$ in the unit ball about $\theta_*$. This intersection is full dimensional, so $\mb P (X_j \in S_i) = c_i > 0$. Let $b_n$ be the probability that there is some $\theta \in \FMset (X_1, \dots, X_n)$ in $B_1(\theta_*)$. By the argument above, we can bound $p_n$ below using $b_n, c_i$:
		\begin{align*}
			p_n &\geq \mb P(A = \{ \A(X_j - \theta), j=1,\dots,n \}) \\
			&\geq \mb P(\{\emptyset \neq \FMset(X_1,\dots,X_n) \cap B_1(\theta_*)\} \cap \{ \forall C_i, \emptyset \neq C_i \cap \{ X_1, \dots, X_n \} \}) \\
			&\geq 1 - (1-b_n) - \textstyle \sum_{i \in [r]} (1-c_i)^n
		\end{align*}
		By the consistency of $\theta_n$ (\Cref{thm:consistent_FM}), $b_n \rightarrow 1$. So this lower bound converges to 1.
		\begin{figure}[ht!]
			\centering
			\begin{tikzpicture}
				\draw (0,-1) -- (1, 0) -- (0,1) -- (-1,0) -- (0,-1);
				\draw (-6,-1) -- (6,-1);
				\draw (-6,1) -- (6,1);
				\draw (-1,-3) -- (-1,3);
				\draw (1,-3) -- (1,3);
				\draw[<->] (-1,0) -- (0,0) node[anchor=north east] {$1$};
				\filldraw (0,0) circle (2pt) node[anchor=north west] {\small$\theta_*$};
				\filldraw (-0.3,0.4) circle (2pt) node[anchor=south east] {\small$\theta_n$};
				\node at (4,2) {$S_2$};
				\node at (-4,2) {$S_1$};
				\node at (4,-2) {$S_3$};
				\node at (-4,-2) {$S_4$};
				\filldraw (2,2) circle (2pt) node[anchor=south west] {\small$X_{j(2)}$};
				\filldraw (-1.8,-2) circle (2pt) node[anchor=north east] {\small$X_{j(4)}$};
				\filldraw (2,-2) circle (2pt) node[anchor=north west] {\small$X_{j(3)}$};
				\filldraw (-2.4,2.3) circle (2pt) node[anchor=south east] {\small$X_{j(1)}$};
			\end{tikzpicture}
			\caption{A visualisation of the cones $S_i$ (defined in the proof of \cref{thm:uniqueness_prob_convergence}) in the case of the $\ell_1$ norm, and a sample such that every cone contains some data point. As each cone contains a data point, whenever there is a Fr\'echet mean within the unit ball of $\theta_*$, it must be a unique Fr\'echet mean.}
			\label{fig:unique_prob_to_1}
		\end{figure}
	\end{proof}
	
	We have shown that the probability of uniqueness converges to one, now we argue that it will never be exactly one. This is best seen by considering a sample close to a  line $\vec a \R$ for some $\vec a \in A$, such that we have $\A(X_{j_1}-X_{j_2}) = \vec a$ or $-\vec a$ for all $j_1 \neq j_2 \leq n$; such a sample will occur with positive probability. For such a sample, the Fr\'echet mean set is a $(k-1)$ dimensional polytope contained in the hyperplane $\{\vec x: \vec a^{\T}\vec x = n^{-1} \sum_{j\leq n} \vec a^{\T} X_j \}$, and so is not unique. \Cref{fig:non-unique_FM_sets} shows an example of this using the $\ell_{\infty}$-norm.
	
	\begin{figure}[ht]
		\centering
		\begin{tikzpicture}[scale=1.8]
			
			\coordinate (A) at (0,0);
			\coordinate (B) at (-1,0.25);
			\coordinate (C) at (2,-0.25);
			\coordinate (D) at (1,0.5);            
			
			\draw ($(A) + 0.5*(1,1)$) -- ($(A) + 0.5*(-1,1)$) -- ($(A) + 0.5*(-1,-1)$) -- ($(A) + 0.5*(1,-1)$) -- ($(A) + 0.5*(1,1)$);
			\draw ($(B) + 1.5*(1,1)$) -- ($(B) + 1.5*(-1,1)$) -- ($(B) + 1.5*(-1,-1)$) -- ($(B) + 1.5*(1,-1)$) -- ($(B) + 1.5*(1,1)$);
			\draw ($(C) + 1.5*(1,1)$) -- ($(C) + 1.5*(-1,1)$) -- ($(C) + 1.5*(-1,-1)$) -- ($(C) + 1.5*(1,-1)$) -- ($(C) + 1.5*(1,1)$);
			\draw ($(D) + 0.5*(1,1)$) -- ($(D) + 0.5*(-1,1)$) -- ($(D) + 0.5*(-1,-1)$) -- ($(D) + 0.5*(1,-1)$) -- ($(D) + 0.5*(1,1)$);
			\filldraw (A) circle (1pt) node[anchor=north] {\small $X_1 = (0,0)$};
			\filldraw (B) circle (1pt) node[anchor=north] {\small $X_3 = (-4,1)$};
			\filldraw (C) circle (1pt) node[anchor=north] {\small $X_4 = (8,-1)$};
			\filldraw (D) circle (1pt)  node[anchor=south] {\small $X_2 = (4,2)$};
			\draw[red, line width = 3pt]  (0.5,0) -- (0.5,0.5);
		\end{tikzpicture}
		\caption{A data configuration which occurs with positive probability and has non-unique $\ell_\infty$ Fr\'echet means.}
		\label{fig:non-unique_FM_sets}
	\end{figure}
	
	\begin{remark}\label{rmk:non_unique_for_2n}
		The probability of uniqueness is always zero for samples of size 2. With probability 1, $X_2$ is in $X_1 + \relint \R_+ F$ for some facet $F$ of $B$. In this case, the Fr\'echet mean set $\FMset$ is a $k-1$ dimensional polytope contained in the hyperplane $\{ \vec x : \vec a^\T\vec x = \vec a^\T(X_1+X_2)/2 \}$ where $\vec a$ is the normal vector of the facet $F$.
	\end{remark}
	
	We've verified that the population Fr\'echet mean is unique for positive density distributions. We have also used the consistency of sample Fr\'echet means to deduce that the probability of a sample having a unique Fr\'echet mean converges to 1 (but never equals 1) as our sample size increases. In the following subsection, we study the uniqueness behaviour of Fr\'echet means for small samples.
	
	\section{Sample Threshold for Unique Fr\'echet Means}
	\label{sec:FM_uniqueness}
	
	
	In this section, we identify the combinatorial conditions that dictate the sample size needed for the sample Fr\'echet mean to be unique with positive probability. We show that this sample size is a threshold in that there is some value $N(B)$ such that the probability of a unique Fr\'echet mean is positive for a sample of size $n$ if and only if $n \geq N(B)$. For the $\ell_1$ norm, this sample threshold is just 3, but for the $\ell_{\infty}$ norm, this sample threshold is equal to the dimension of the space, $k$.
	
	To understand the dimension of Fr\'echet mean sets, we study which faces of balls centred about our data points can be intersected to define the Fr\'echet mean set (\cref{cor:FMisPolytope}). The dimension of the Fr\'echet mean set $\FMset$ is identifiable from the minimal faces whose intersection is the Fr\'echet mean set. We call these minimal faces the \emph{face type} of a sample. We show that the face type of a sample is almost surely well-defined and determines the dimension of $\FMset$, then we identify which face types occur with positive probability by considering subgradients. To define the face type of a sample, we will assume that no data point is also a Fr\'echet mean of the sample, which is true almost surely.
	
	\begin{lemma}\label{lemma:data_is_not_FM}
		If $\esssupp(\mathbb P) = \R^k$ then with probability one, the random sample points $ \sample \overset{i.i.d.}{\sim} \mathbb P$ are not Fr\'echet means of the full sample; for all $i$, $X_i \notin \FMset(\sample)$.
	\end{lemma}
	\begin{proof}
		We claim $X_1$ is a Fr\'echet mean of $\sample$ iff $X_1$ is a Fr\'echet mean of $X_2, \dots, X_n$. Denote the Fr\'echet function for the sample $\sample$ as $f(\vec x)$, and the Fr\'echet function for $X_2, \dots, X_n$ as $f_{\hat{1}}(\vec x)$, so:
		\[
		f(\vec x) = n^{-1} \sum_{i=1}^n d(X_i, \vec x), \qquad
		f_{\hat{1}}(\vec x) = (n-1)^{-1} \sum_{i=2}^n d(X_i, \vec x).
		\]
		The subgradient of $d(X_1, \cdot)^2$ at $X_1$ is 0, so by Theorem 23.8 in \cite{rockafellar1970convex}:
		\[
		n \partial f(X_1) = (n-1) \partial f_{\hat{1}}(X_1).
		\]
		The point $X_1$ is a Fr\'echet mean of $\sample$ if and only if $0 \in \partial f(X_1)$, and is a Fr\'echet mean of $X_2, \dots X_{n}$ if and only if $0 \in \partial f_{\hat{1}}(X_1)$. These conditions are equivalent, so our claim is proved.
		
		We now show that the sample points $\sample$ are almost never Fr\'echet means of the full sample $\sample$. Let $A_1$ be the event that $X_1$ is a Fr\'echet mean of the full sample, or equivalently $X_1 \in \FMset(X_2, \dots, X_n)$. By Corollary \ref{cor:FMisPolytope}, the set $\FMset(X_2,\dots, X_n)$ has co-dimension $\geq 1$, and hence has $\mathbb P$-measure zero. Therefore $\mathbb P(A_1) = 0$. We can define $A_2, \dots, A_n$ similarly, and by taking a union over $A_1, \dots, A_n$ we conclude that the Fr\'echet mean set intersects the sample $\sample$ with probability zero.
	\end{proof}
	
	\subsection{Face Types}
	
	As shown in \Cref{cor:FMisPolytope}, the Fr\'echet mean set is an intersection of faces of balls centred at each data point. We will define the face type of a sample as the minimal sequence of faces of balls where the Fr\'echet mean set is their intersection. The lemma below ensures this is almost surely well-defined.
	
	\begin{lemma}
		\label{lem:unique_least_face_seq}
		Let $\sample$ be a finite sample with Fr\'echet mean set $\FMset$ such that $X_i \notin \FMset$ for every $i$. Let $d_i = d_B(X_i, \FMset)$. Then the set of tuples of faces
		\[
		S_F(\sample) = \bigg\{ (F'_1, \dots, F'_n) \ : \textstyle \bigcap_{i = 1}^n (X_i - d_i F'_i) = \FMset \bigg\} \subseteq \mathfrak{F}(B)^n
		\]
		has a least element $(F_1, \dots, F_n) \in S_F$ with respect to coordinate-wise containment.
	\end{lemma}
	
	\begin{proof}
		By \Cref{cor:FMisPolytope}, $S_F(\sample)$ is not empty. Define $(F_1, \dots, F_n)$ by $F_i = \bigcap_{(F'_1, \dots, F'_n) \in S_F} F'_i$.
		Then for all $(F'_1, \dots, F'_n) \in S_F(\sample)$ we have that $(F_1, \dots, F_n) \leq (F'_1, \dots, F'_n)$. \\
		It suffices to show that $(F_1, \dots, F_n) \in S_F(\sample)$. This follows from:
		\begin{align*}
			\textstyle \bigcap_{i = 1}^n (X_i - d_i F_i) &= \textstyle \bigcap_{i = 1}^n \bigg(X_i - d_i \big( \textstyle \bigcap_{(F'_1, \dots, F'_n) \in S_F} F'_i\big)\bigg) \\
			%&= \bigcap_{i = 1}^n \bigcap_{(F'_1, \dots, F'_n) \in S_F} (X_i - d_i F'_i) \\
			&= \textstyle \bigcap_{(F'_1, \dots, F'_n) \in S_F} \bigcap_{i = 1}^n (X_i - d_i F'_i) \\
			&= \textstyle \bigcap_{(F'_1, \dots, F'_n) \in S_F} \FMset = \FMset.
		\end{align*}
	\end{proof}
	Using the lemma above, we can define the face type of a sample:
	\begin{definition}[Face Type]
		Consider a sample $\sample$ such that no data point is a Fr\'echet mean of the full dataset. Such a sample is defined to have face type $(F_1, \dots, F_n)$ about the Fr\'echet mean set $\FMset$ when $(F_1, \dots, F_n)$ is the least element of $S_F(\sample)$.
	\end{definition}
	
	By \cref{lemma:data_is_not_FM} and \cref{lem:unique_least_face_seq}, the face type of a sample is well-defined almost surely. However, we will use a more geometric description of face type through the rest of this section. Both here and throughout the rest of this section we use $d_i \coloneqq d_B(X_i, \FMset)$ which is well-defined by \cref{cor:FMisPolytope}.
	
	\begin{lemma} \label{lemma:geometry_comb_type}
		Let $\sample$ be a sample with face type $(F_1, \dots F_n)$ and $d_i = d_B(X_i,\FMset)$. Then we have that 
		$
		\relint \FMset = \textstyle\bigcap_{i = 1}^n \relint (X_i - d_i F_i).
		$
		A sample has face type $(F_1, \dots, F_n)$ if and only if there is some $\theta \in \relint \FMset$ such that for all $i \in [n]$, $X_i - \theta \in \relint (\R_+F_i)$.
	\end{lemma}
	\begin{proof}
		By definition, $\FMset \subseteq (X_1 - d_1 F_1)$.     Suppose that $\relint \FMset$ does not intersect $\relint (X_1 - d_1 F_1)$, so $\relint \FMset \subseteq (X_1 - d_1 \partial F_1)$. By Corollary \ref{cor:convex_disjoint_from_relint}, there is some face $F_1'$ of $F_1$ satisfying $\FMset \subseteq X_1 - d_1 F_1'$. Hence $(F_1', F_2, \dots, F_n) \in S_F(X_1, \dots,X_n)$, contradicting the definition of $(F_1, \dots, F_n)$. Therefore (for any $i$, by symmetry) $\relint \FMset$ intersects $\relint (X_i - d_i F_i)$, and by \cite[Theorem 6.5]{rockafellar1970convex}:
		\begin{align*}
			\relint (\FMset \cap (X_i - d_i F_i)) &= \relint \FMset \cap \relint (X_i - d_i F_i), \\
			\relint \FMset &\subseteq \relint (X_i - d_i F_i).
		\end{align*}
		Applying \cite[Theorem 6.5]{rockafellar1970convex} again to $X_1 - d_1 F_1, \dots, X_n - d_n F_n$:
		\begin{align*}
			\relint \FMset = \textstyle \bigcap_{i = 1}^n \relint (X_i - d_i F_i).
		\end{align*}
		Suppose our sample has face type $F_1, \dots, F_n$. Then for all $\theta \in \relint \FMset$, we have $\theta \in \relint (X_i - d_i F_i)$ and therefore $X_i - \theta \in \relint \R_+ F_i$.
		
		To show the converse, consider $\theta \in \relint \FMset$ such that for all $i$, $X_i - \theta \in \relint (\R_+ F_i)$. We define $d_i = d_B(X_i, \theta)$. Assume that $(\sample)$ has the face type $F_1', \dots, F_n'$. Then by the argument above, $\theta \in \relint(X_i - d_i F_i')$, implying that $X_i - \theta \in d_i \relint F_i' \subseteq \relint (\R_+ F_i')$. The fans $\{ \relint (\R_+F) : F \in \mathscr F (B) \}$ partition $\R^k \setminus \{ \vec 0 \}$, so it must be the case that $F_i' = F_i$ and our sample has type $(F_1, \dots, F_n)$. 
	\end{proof}
	
	As the relative interior of the Fr\'echet mean set is the intersection of the relative interiors of the faces $X_i - d_i F_i$, we can deduce the dimension of $\FMset$ from the face type of the sample.
	
	\begin{lemma}\label{lemma:affine_intersection}
		If $W_1,\ldots,W_m$ are relatively open and $\textstyle \bigcap_{i=1}^m W_i \neq \emptyset$, then
		$$
		\textstyle\bigcap_{i=1}^m\mathrm{aff}(W_i) = \mathrm{aff} \left( \textstyle\bigcap_{i=1}^m W_i \right).
		$$
	\end{lemma}
	\begin{proof}
		It suffices to consider the $m=2$ case. The space $\text{aff}(W_1) \cap \text{aff}(W_2)$ is affine and contains $W_1 \cap W_2$, so we get that $\text{aff}(W_1 \cap W_2) \subseteq \text{aff}(W_1) \cap \text{aff}(W_2)$.\\
		Now fix $\vec y \in \text{aff}(W_1) \cap \text{aff}(W_2)$, $\vec x \in W_1 \cap W_2$. As $W_1,W_2$ are relatively open, there exists some $\epsilon > 0$ such that $(1-\epsilon)\vec x + \epsilon \vec y \in W_1 \cap W_2$. Hence $\vec y \in \text{aff}(\{ \vec x, (1-\epsilon)\vec x + \epsilon \vec y \}) \subseteq \text{aff}(W_1 \cap W_2)$.     
	\end{proof}
	
	\begin{remark}\label{rmk:positive_dist_to_data}
		The face type of a sample is undefined if and only if some data point $X_i$ is a Fr\'echet mean of the full dataset. \Cref{lemma:FM_dists} tells us the Fr\'echet mean is the unique point $X_i$ in this case, while \Cref{lemma:data_is_not_FM} tells us this is a probability zero event. For the rest of this section we exclude this case, assuming that $d_i = d_B(X_i, \FMset) > 0$.
	\end{remark}
	
	\subsection{Type Events}
	
	By \Cref{lemma:affine_intersection}, we can deduce the dimension of a Fr\'echet mean set from the face type of a sample. We must now identify which face types occur with positive probability, and identify which face types correspond to a unique Fr\'echet mean. To determine which face types occur with positive probability, we define the (lifted) type event. Through this section, we use $C_i$ and $G_i$ to denote $\R_+ F_i$ and $F_i^{\Diamond}$ respectively.
	
	\begin{definition}[(Lifted) Type Event]
		Consider the random vector $(\sample) \in \R^{nk}$. In this space, the \emph{type event} $U_{\CType}$ is given by:
		\[
		U_{F_1, \dots, F_n} = \{ (\sample) \in \R^{nk}: (\sample) \text{ has face type }(F_1, \dots F_n) \}.
		\]
		Then we define the \emph{lifted type event} as the set
		\[
		U_{F_1, \dots, F_n, \Theta} = \{ (\sample, \theta) \in \R^{k(n+1)} : \theta \in \relint \FMset, 
		\; X_i - \theta \in \relint C_i \},
		\]
		where $C_i$ denotes the cone $\R_+ F_i$.
	\end{definition}
	
	Note that by \Cref{lemma:geometry_comb_type}, the type event is the image of the lifted type event under the projection $\proj: \R^{k(n+1)} \rightarrow \R^{kn}$ given by $(\sample, \theta) \mapsto (\sample)$.
	
	Our goal is to classify the face types $(\CType)$ for which $U_{\CType}$ has positive Lebesgue measure. Then as $(\sample)$ has a positive density, such a face type occurs with positive probability. It is difficult to calculate the measure of the type events directly, so we will translate this condition into a linear relation on subgradients. For this, we define the \emph{zero subgradient event}.
	
	\begin{definition}[Zero Subgradient Event]
		The zero subgradient event is given by:
		\[
		V_{\CType, \Theta} = \left\{ (\sample, \theta) \in \R^{k(n+1)} : \vec 0 \in \relint \left( \textstyle \sum_{i=1}^n \|X_i-\theta\| G_i \right),\; X_i - \theta \in \relint C_i \right\},
		\]
		where $G_i$ denotes the polar face $F_i^{\Diamond}$, and $C_i$ denotes the cone $\R_+F_i$.
	\end{definition}
	
	Note that when we have $X_i - \theta \in \relint C_i$ for all $i$, the subgradient of the Fr\'echet function is given by $\partial f(\theta) = \sum_{i=1}^n \| X_i - \theta \|G_i$. Hence, by the definition of subgradients, $\theta \in \FMset$ if and only if $\vec 0 \in \sum_{i=1}^n \| X_i - \theta \|G_i$ (\cref{rmk:subgradient}). From this, we might hope that the lifted type event and the zero subgradient event are equivalent, but this is not true. However, they do coincide almost everywhere.
	
	\begin{restatable}{proposition}{EventEquivalence}\label{prop:lifted_type_event}
		The zero subgradient event $V_{\CType, \Theta}$ is a relatively open polyhedron; it is an intersection of finitely many hyperplanes and open half spaces. Furthermore, it coincides with the lifted type event everywhere other than on its boundary:
		\begin{align*}
			\relint V_{\CType, \Theta} &= \relint U_{\CType, \Theta},\\
			\closure V_{\CType, \Theta} &= \closure U_{\CType, \Theta}.
		\end{align*}
	\end{restatable}
	
	The proof of this proposition can be found in \cref{appsec:prop_proof}.
	
	\subsection{Polytope Conditions for Unique Fr\'echet Means}
	
	We can now state and prove the main result of this section (\cref{thm:wwp_dim_condition}), which establishes the conditions for $U_{\CType}$ to occur with positive probability and to produce a unique Fr\'echet mean. We do this by computing the dimension of the polyhedron $V_{\CType,\Theta}$, which then allows us to compute the dimension of $\proj V_{\CType,\Theta}$ and $U_{\CType}$. This result is crucial for not only computing the sample threshold for uniqueness, but also for proving that the sample threshold for uniqueness is well-defined.
	
	\begin{theorem}\label{thm:wwp_dim_condition}
		The type event $U_{\CType}$ and $\proj V_{\CType,\Theta}$ satisfy
		$
		\L_{}(U_{\CType}) = \L(\proj V_{\CType,\Theta}).
		$ \\
		Therefore $U_{\CType}$ has positive probability if and only if
		\begin{gather}
			\vec 0 \in \relint \left( \conv \left(\textstyle \bigcup_{i \in [n]} G_i \right) \right) \label{eq:wpp_cond_possible}, \\
			\text{and} \qquad \textstyle \sum_{i \in [n]} \dim G_i = \dim \textstyle \sum_{i \in [n]} \aff G_i. \label{eq:wpp_cond_probable}
		\end{gather}
		The Fr\'echet mean of a sample in $U_{\CType}$ is unique if and only if 
		\begin{align}\label{eq:wpp_cond_unique}
			\dim \conv \left( \textstyle \bigcup_{i \in [n]} G_i \right) = k.
		\end{align}
	\end{theorem}
	
	To understand this result, we return to one of our examples from \cref{fig:FM_sets}, displayed again in \cref{fig:Face_Type_Example}, which shows a sample with face type $(F_1,F_2,F_3) = (\vec e_1+\vec e_2, \conv \{ -\vec e_2 \pm \vec e_1 \}, \conv \{ -\vec e_1 \pm \vec e_2 \})$. We can test whether this face type satisfies the three conditions in \cref{thm:wwp_dim_condition}; for this face type, note that $(G_1, G_2, G_3) = (\conv \{\vec e_1, \vec e_2 \}, -\vec e_2, -\vec e_1)$. \\
	The first condition \cref{eq:wpp_cond_possible} verifies whether a sample can have the face type $(F_1,F_2, F_3)$. This is shown to be true in \cref{fig:Face_Type_Example}, and in fact $\conv \left(\textstyle \bigcup_{i \in [n]} G_i \right) = \conv \{\pm\vec e_1,\pm\vec e_2 \}$ which contains $\vec 0$ in its relative interior.\\
	The second condition \cref{eq:wpp_cond_probable} verifies whether a sample has face type $(F_1,F_2,F_3)$ with positive probability. The dimensions of $G_1,G_2,G_3$ are 1,0,0 respectively, while $\sum_{i \leq 3} \aff G_i = \aff \{ -\vec e_1, -\vec e_2 \}$, which has dimension $1=1+0+0$. This verifies that such a face type can occur with positive probability; this configuration of data points around a Fr\'echet mean is generic.\\
	Finally, the third condition \cref{eq:wpp_cond_unique} verifies whether having face type $(\CType)$ implies that our sample has a unique Fr\'echet mean. We know that $\conv \left(\textstyle \bigcup_{i \in [n]} G_i \right) = \conv \{\pm\vec e_1,\pm\vec e_2 \}$ which has dimension 2, and therefore any sample with this face type will have unique Fr\'echet mean. \\
	
	\begin{figure}
		\centering
		\begin{tikzpicture}[scale=0.9]
			\draw (0,0) -- (-4, 0) -- (-4,-4) -- (0,-4) -- (0,0);
			\draw (-0.5,0) -- (1.5, 0) -- (1.5,2) -- (-0.5,2) -- (-0.5,0);
			\draw (0,0.5) -- (0,-1.5) -- (2,-1.5) -- (2,0.5) -- (0,0.5);
			\filldraw (-2,-2) circle (1pt) node[anchor=north east] {\small$\vec x_1 = (-4,-4)$};
			\filldraw (0.5,1) circle (1pt)  node[anchor=south] {\small$\vec x_2 = (1,2)$};
			\filldraw (1,-0.5) circle (1pt)  node[anchor=north west] {\small$\vec x_3 = (2,-1)$};
			\filldraw[red] (0,0) circle (3pt);
		\end{tikzpicture}  
		\caption{A sample with face type $(\CType) = (\vec e_1+\vec e_2, \conv \{ - \vec e_2 \pm \vec e_1 \}, \conv \{ - \vec e_1 \pm \vec e_2 \})$, and its unique Fr\'echet mean highlighted in red.}
		\label{fig:Face_Type_Example}
	\end{figure}
	
	We prove \cref{thm:wwp_dim_condition} via the following three lemmas.
	
	\begin{lemma}\label{lemma:possible}
		The zero subgradient event $V_{\CType,\Theta}$ is non-empty if and only if $\vec 0 \in \relint \left( \conv \left(\textstyle \bigcup_{i \in [n]} G_i \right) \right)$.
	\end{lemma}
	
	\begin{proof}
		First, we have shown in the proof of \cref{prop:lifted_type_event} that
		\begin{equation}
			V_{\CType, \Theta} = L_{\CType}^{-1}(W) \textstyle \bigcap \bigcap_{i \in [n]} \{ (\sample, \theta) : X_i - \theta \in \relint C_i \},\label{eq:zero_subgrad_event_as_intersection}
		\end{equation}
		where $W$ is the relatively open polyhedral cone $\{\vec w \in \R^n_{+} : \vec 0 \in \sum_{i \in [n]} w_i \relint G_i\}$ (\cref{lemma:polyhedral_cone_W}). 
		
		We will show that $V_{\CType,\Theta}$ is not empty if and only if $W$ is non-empty. Suppose $(X_1,\dots, X_n,\theta) \in V_{\CType,\Theta}$; then by \cref{eq:zero_subgrad_event_as_intersection}, $W$ is non-empty. Conversely, suppose $\vec w \in W$. Fix an arbitrary $\theta$ and pick some $X_i \in \theta + w_i \relint F_i$. By construction, $X_i - \theta \in \relint C_i$ and $\| X_i - \theta \|_B = w_i$. By the definition of $W$ we have $
		\vec 0 \in \relint \left( \textstyle \sum_{i \in [n]} \| X_i - \theta \|_B G_i \right),
		$ so $(\sample,\theta) \in V_{\CType,\Theta}$.
		
		Next, note that \cite[Theorem 6.9]{rockafellar1970convex} states that for non-empty convex $G_1, \dots, G_n \subset \R^k$, then
		\begin{equation}\label{eq:rockafellar_thm6.9}
			\relint \left( \conv \left( \textstyle \bigcup_{i \in [n]} G_i \right) \right) = \bigcup \left \{ \textstyle \sum_{i \in [n]} \lambda_i \relint G_i : \lambda_i > 0, \sum_{i \in [n]} \lambda_i = 1 \right \}.
		\end{equation}
		From this we see that $\vec 0 \in \relint \left( \conv \left(\textstyle \bigcup_{i \in [n]} G_i \right) \right)$ if and only if $W$ is non-empty. We conclude that $V_{\CType, \Theta}$ and $V_{\CType}$ are non-empty if and only if $\vec 0 \in \relint \left( \conv \left(\textstyle \bigcup_{i \in [n]} G_i \right) \right)$.
	\end{proof}
	
	\begin{lemma}\label{lemma:probable}
		The type event $U_{\CType}$ has positive probability if and only if the polyhedron $\proj V_{\CType,\Theta}$ has dimension $nk$.
	\end{lemma}
	
	\begin{proof}
		The zero subgradient event $V_{\CType,\Theta}$ is Borel measurable as it is a relatively open polyhedron. Consequently, $\proj V_{\CType,\Theta}$ is also a relatively open polyhedron and so also Borel measurable.
		We also have by Proposition \ref{prop:lifted_type_event} and Theorem 6.6 of \cite{rockafellar1970convex}:
		\begin{align*}
			\relint U_{\CType} &= \relint (\proj U_{\CType,\Theta})\\
			&= \proj (\relint U_{\CType,\Theta}) \\
			&= \proj (\relint V_{\CType,\Theta}) \\
			&= \relint (\proj V_{\CType,\Theta}),\\
			U_{\CType} &= \proj U_{\CType,\Theta} \\
			&\subseteq \proj (\closure U_{\CType,\Theta}) \\
			&= \proj (\closure V_{\CType,\Theta}) \\
			&\subseteq \closure (\proj V_{\CType,\Theta}).
		\end{align*}
		Hence $\proj V_{\CType,\Theta}$ is a measurable set which differs from $U_{\CType}$ only on its boundary. The boundary of $\proj V_{\CType,\Theta}$ is a measure zero set as it is a relatively open polyhedron, so by the completeness of the Lebesgue measure, $U_{\CType}$ is measurable and $\L(U_{\CType}) = \L(\proj V_{\CType,\Theta})$.
		
		As a polyhedron, $\proj V_{\CType,\Theta}$ has positive Lebesgue measure if and only if it has full dimension $nk$. 
	\end{proof}
	
	
	\begin{lemma} \label{lemma:unique}
		A sample with face type $(\CType)$ will have a Fr\'echet mean set of dimension $$k - \dim \conv \left(\textstyle \bigcup_{i \in [n]} G_i \right).$$ 
	\end{lemma}
	
	\begin{proof}
		We recall \cref{lemma:affine_intersection}; we have that $\relint \FMset = \bigcap_{i \in [n]} \relint (X_i - d_i F_i)$, so
		\begin{align}
			d  &= \dim \textstyle\bigcap_{i \in [n]} \aff (X_i - d_i F_i), \nonumber\\
			&=\dim \left\{ \theta + \vec v : \theta \in \textstyle\bigcap_{i \in [n]} (X_i - d_i F_i), \vec v \in G_i^{\perp} \text{ for all }i \in [n] \right\}, \nonumber\\
			&=\dim \left(\textstyle \bigcup_{i \in [n]} G_i \right)^{\perp}, \nonumber\\
			&=k - \dim \conv \left(\textstyle \bigcup_{i \in [n]} G_i \right). \label{eq:expr_for_d}
		\end{align}
	\end{proof}
	
	We can now prove \cref{thm:wwp_dim_condition}.
	
	\begin{proof}
		By \cref{lemma:possible} and $\cref{prop:lifted_type_event}$, if \cref{eq:wpp_cond_possible} does not hold, then $V_{\CType,\Theta}$ and $U_{\CType}$ are empty. We therefore assume that \cref{eq:wpp_cond_possible} holds such that $V_{\CType,\Theta}$ is non-empty, and then compute its dimension. 
		
		As $V_{\CType,\Theta}$ is non-empty, we can use \cref{lemma:affine_intersection}:
		\begin{align*}
			\text{aff}(V_{\CType,\Theta}) &= \aff L_{\CType}^{-1}(W) \textstyle \bigcap \bigcap_{i \in [n]} \text{aff} \{ (\sample, \theta) : X_i - \theta \in \relint C_i \}.
		\end{align*}
		Observe that
		\begin{align*}
			\aff L_{\CType}^{-1}(W) &= \{\textstyle\sum_{i \in [n]} \lambda_i \vec x_i : L_{\CType} \vec x_i \in W, \lambda_i \in \R, \sum_{i \in [n]} \lambda_i = 1\}, \\
			&= \{L_{\CType}^{-1}\textstyle\sum_{i \in [n]} \lambda_i \vec w_i : \vec w_i \in W, \lambda_i \in \R, \sum_{i \in [n]} \lambda_i = 1\}, \\
			&= L_{\CType}^{-1} \aff(W).
		\end{align*}
		And therefore 
		\begin{align*}
			\text{aff}(V_{\CType,\Theta}) &= L_{\CType}^{-1}(\aff (W)) \textstyle \bigcap \bigcap_{i \in [n]} \text{aff} \{ (\sample, \theta) : X_i - \theta \in \relint C_i \}.
		\end{align*}
		Recall that $L_{\CType}$ is defined as $L_{\CType}: (\sample,\theta) \mapsto \left(\vec a_1^{\T} (X_1-\theta), \dots, \vec a_n^{\T} (X_n-\theta)\right)$, where $\vec a_i$ is some element of $\A(F_i)$. The affine hull of $C_i$ contains those points $\vec x$ for which $\vec a^{\T}\vec x$ is constant for all $\vec a \in \A(F_i)$; it is the orthogonal to $\aff \A(F_i) = \aff G_i$. Hence $\aff C_i$ has dimension 
		\begin{align*}
			\dim \aff C_i &= k - \dim \aff \A(F_i), \\
			&= k - \dim G_i.
		\end{align*}
		Noting this, we can pick $\vec a_i^1 \dots, \vec a_i^{\dim G_i}$ in $\A(F_i)$ such that $\vec a_i, \vec a_i^1 \dots, \vec a_i^{\dim G_i}$ forms an affine basis of $\aff \A(F_i)$, and $\aff C_i$ contains those $\vec x$ such that for all $j \in [\dim G_i]$, we have $\vec a_i^{j\T}\vec x = \vec a_i^{\T}\vec x$. Then 
		\begin{align*}
			\aff \{ (\sample,\theta) : X_i - \theta \in C_i \} &= \{ (\sample,\theta) : X_i - \theta \in \aff C_i \}, \\
			&= \{ (\sample,\theta) : \forall \, j \in [\dim G_i], 
			\; (\vec a_i^j - \vec a_i)^{\T}(X_i - \theta) = 0 \}.
		\end{align*}
		We define $M_i$ to be the linear map sending $(\sample,\theta)$ to $\left((\vec a_i^{j}-\vec a_i)^{\T}(X_i-\theta)\right)_{j \in [\dim G_i]}$, such that $\aff \{ (\sample,\theta) : X_i - \theta \in \relint C_i \} = M_i^{-1}(\vec 0)$. Then 
		\begin{align*}
			\text{aff}(V_{\CType,\Theta}) &= L_{\CType}^{-1}(\aff(W)) \cap M_1^{-1}(\vec 0) \cap \dots \cap M_n^{-1}(\vec 0), \\
			&= (L_{\CType},M_1,\dots,M_n)^{-1}(\aff(W),\vec 0,\dots,\vec 0).
		\end{align*}
		The linear map $(L_{\CType},M_1,\dots,M_n) : \R^{k(n+1)} \rightarrow \R^{n + \sum_i \dim G_i}$ is given by the matrix:
		$$
		M = \begin{pmatrix}
			\vec a_1^{\T} & \dots & \vec 0^\T & -\vec a_1^{\T} \\
			\vdots & \ddots & \vdots & \vdots \\
			\vec 0^\T & \hdots & \vec a_n^{\T} & -\vec a_n^{\T} \\
			(\vec a_1^1 - \vec a_1)^{\T}& \hdots & \vec 0^\T & (\vec a_1 - \vec a_1^1)^{\T} \\
			\vdots & \vdots & \vdots & \vdots \\
			(\vec a_1^{\dim G_1}- \vec a_1)^\T& \hdots & \vec 0^\T & (\vec a_1^{\dim G_1} - \vec a_1)^\T \\
			\vdots & \ddots & \vdots & \vdots \\
			\vec 0^\T & \hdots & (\vec a_n^1 - \vec a_n)^\T & (\vec a_n - \vec a_n^1)^{\T} \\
			\vdots & \vdots & \vdots & \vdots \\
			\vec 0^\T & \hdots & (\vec a_n^{\dim G_n} - \vec a_n)^{\T} & (\vec a_n - \vec a_n^{\dim G_n})^{\T} \\
		\end{pmatrix}.
		$$
		Subtracting row $i$ from the rows corresponding to $M_i$, we see the rank of this matrix is given by:
		\begin{align*}
			\textstyle \sum_{i \in [n]} \dim \text{span} (\vec a_i, \vec a_i^1, \dots, \vec a_i^{\dim G_i}) &= \textstyle \sum_{i \in [n]} 1+ \dim \aff (\vec a_i, \vec a_i^1, \dots, \vec a_i^{\dim G_i}), \\
			&= \textstyle \sum_{i \in [n]} 1+ \dim G_i, \\
			&= n + \textstyle \sum_{i \in [n]} \dim G_i.
		\end{align*}
		Where the first line is because $\aff (\vec a_i, \vec a_i^1, \dots, \vec a_i^{\dim G_i}) = \aff \A(F_i)$ doesn't contain the origin; for all $\vec a \in \aff \A(F_i), \, \vec a^\T (X_i-\theta) = \|X_i - \theta\| > 0$.\\    
		It follows that the matrix $M$ is full rank.  Let $T: \R^{n + \sum_i \dim G_i} \rightarrow \R^{n + \sum_i \dim G_i-\dim W}$ be the linear projection onto the orthogonal complement of $(\aff(W),\vec 0)$ such that the kernel of $T$ is exactly $(\aff(W),0)$. 
		Then by construction, the map $T \circ M: \R^{k(n+1)} \rightarrow \R^{n + \sum_i \dim G_i-\dim W}$ is surjective with kernel $\aff(V_{\CType,\Theta})$ and by the rank-nullity theorem:
		\begin{align*}
			\dim V_{\CType,\Theta} &= k(n+1) - \rank (T \circ M) = k(n+1) - n - \textstyle \sum_{i \in [n]} \dim G_i + \dim W.
		\end{align*}
		Next, we compute the dimension of $U_{\CType}$. For this, we will use \cite[Theorem 2.4]{balas1998dimension}, which states:
		\begin{theorem*}[Theorem 2.4 of \cite{balas1998dimension}]
			Consider a polytope $Q \subset \R^p \times \R^q$ with affine hull is given by $\{ (\vec u,\vec x) \in \R^p \times \R^q : A\vec u+B\vec x = \vec b\}$. Then $\dim \proj_{\R^q}(Q) = \dim Q - \ker(A)$.
		\end{theorem*}
		We apply this theorem to $V_{\CType,\Theta}$; define matrices $M^1, M^2$ and vector $\vec b$ to be such that
		\[
		\aff V_{\CType,\Theta} = \{(\vec x_1, \dots, \vec x_n,\theta) \in \R^{nk} \times \R^{k} :  M^1 (\vec x_1, \dots, \vec x_n) + M^2 \theta = \vec b \}.
		\]
		Note that for $(X_1, \dots, X_n) \in U_{\CType}$, the set $\{\theta \in \R^k : M^2 \theta = \vec b - M^1(\sample) \}$ is now the affine hull of the Fr\'echet mean set $\aff \FMset(\sample)$. Therefore $\Ker M^2 = d$, where $d$ is the dimension of the Fr\'echet mean set for a sample with face type $\CType$. Using \cite[Theorem 2.4]{balas1998dimension}, we see that
		\begin{align}
			\dim U_{\CType} &= \dim \proj V_{\CType,\Theta}, \nonumber\\
			&= \dim V_{\CType,\Theta} - d, \nonumber\\
			&= nk + k - n - \textstyle\sum_{i \in [n]} \dim G_i + \dim W - d.\label{eq:init_U_dim}
		\end{align}
		We now look to express $\dim W$ in terms of the polytope $B$. Recall that we have assumed $W$ to be non-empty. We will express the dimension of $W$ in terms of the \emph{Cayley polytope} of $G_1, \dots, G_n$ which we denote by $P_{G_1, \dots, G_n}$. This is given by
		\begin{align*}
			P_{G_1,\dots,G_n} &= \conv \{ \vec e_i \times G_i \subset \R^n \times \R^k : i \in [n] \}, \\
			&= \left \{ (\vec w, \vec g) \in \R^n \times \R^k: \vec w \in \Delta_{n-1}, \, \vec g = \textstyle \sum_{i \in [n]} w_i \vec g_i, \, \vec g_i \in G_i \right \},
		\end{align*}
		where $\Delta_{n-1}$ denotes the $n-1$ dimensional simplex. By examining this definition, we see that $\vec w \in W \bigcap \Delta_{n-1}$ if and only if $(\vec w, \vec 0) \in \relint P_{G_1, \dots, G_n}$. As $W$ is a cone in the positive orthant, $W$ being non-empty is equivalent to $W \bigcap \Delta_{n-1} = \relint P_{G_1, \dots, G_n} \bigcap (\R^n \times \vec 0)$ being non-empty.
		
		We once again use \cite[Theorem 2.4]{balas1998dimension}; we now take $M^1,M^2, \vec b$ to be such that 
		$$
		\aff P_{G_1, \dots, G_n} = \{(\vec w,\vec g) \in \R^n \times \R^k : M^1 \vec w + M^2 \vec g = \vec b \}.
		$$
		Note that $\vec b - M^2 \vec 0$ is in the image of $M^1$ as $\relint P_{G_1, \dots, G_n}$ intersects $(\R^n \times \vec 0)$, so by \cite[Theorem 2.4]{balas1998dimension}: 
		\begin{align}
			\dim P_{G_1,\dots,G_n} - \dim \proj_{\R^k}(P_{G_1,\dots,G_n}) &= \ker(M^1), \nonumber\\
			&= \dim \aff \{ \vec w \in \R^n : (\vec w,\vec 0) \in \relint P_{G_1, \dots, G_n} \}, \nonumber\\
			&= \dim \left( \relint P_{G_1, \dots, G_n} \textstyle\bigcap (\R^n \times \vec 0) \right) = \dim \left(W \textstyle\bigcap \Delta_{n-1}\right).\label{eq:cayley_proj_dim}
		\end{align}
		Similarly, for any $\vec w \in \relint \Delta_{n-1}$, we have that $\vec b - M^1 \vec w$ is in the image of $M^2$. We pick $\vec w = (1/n,\dots, 1/n)$:
		\begin{align*}
			\dim P_{G_1,\dots,G_n} &= \dim \proj_{\R^n}(P_{G_1,\dots,G_n})+ \ker(M^2) \\
			&= \dim \Delta_{n-1} + \dim \aff \{ \vec g \in \R^k : ((1/n,\dots, 1/n),\vec g) \in \relint P_{G_1,\dots,G_n} \} \\
			&= n-1+ \dim \aff \left\{ 1/n \textstyle\sum_{i \in [n]} \vec g_i: \vec g_i \in \relint G_i \right\} \\
			&= n-1+ \dim \textstyle \sum_{i \in [n]} \aff G_i.
		\end{align*}
		We can write out $\dim \proj_{\R^k}(P_{G_1,\dots,G_n})$ explicitly:
		\begin{align*}
			\dim \proj_{\R^k}(P_{G_1,\dots,G_n}) &= \dim \left\{ \textstyle \sum_{i \in [n]} w_i \vec g_i \in \R^k : \vec g_i \in G_i, \vec w \in \Delta_{n-1} \right\}, \\
			&= \dim \conv \left( \textstyle \bigcup_{i \in [n]} G_i \right).
		\end{align*}
		Combining the computations above with \cref{eq:cayley_proj_dim}, we see that
		\begin{align*}
			\dim W &= 1 + \dim (W \textstyle\bigcap \Delta_{n-1}), \\
			&= 1 + \dim P_{G_1,\dots,G_n} - \dim \proj_{\R^k}(P_{G_1,\dots,G_n}), \\
			&= n+\dim \textstyle \sum_{i \in [n]} \aff G_i - \dim \conv \left( \textstyle \bigcup_{i \in [n]} G_i \right).
		\end{align*}
		Finally, using \cref{lemma:unique}, we have that $d = k - \dim \conv \left(\textstyle \bigcup_{i \in [n]} G_i \right)$.
		Inputting our expressions for $d$ and $\dim W$ into \cref{eq:init_U_dim}:
		\begin{align*}
			\dim U_{\CType} &= nk + k - n - \textstyle \sum_{i \in [n]} \dim G_i + \dim W - d, \\
			&= nk - \textstyle \sum_{i \in [n]} \dim G_i + \dim \sum_{i \in [n]} \aff G_i.
		\end{align*}
		Therefore $\L(U_{\CType}) > 0$ if and only if $\textstyle \sum_{i \in [n]} \dim G_i = \dim \sum_{i \in [n]} \aff G_i$ as required. \\
		The final claim of \cref{thm:wwp_dim_condition} follows from \cref{eq:expr_for_d}; $d=0 \Leftrightarrow k = \dim \conv \left( \bigcup_{i \in [n]} G_i \right)$.
	\end{proof}
	
	\subsection{Defining Unique Fr\'echet Mean Sample Thresholds}
	
	\cref{thm:wwp_dim_condition} is the primary theoretical result of this section, explicitly stating the polytope conditions that dictate when a finite sample can have a unique Fr\'echet mean. The question of whether a sample of size $n$ has a positive probability of producing a unique Fr\'echet mean now equates to whether there is a face type $(\CType)$ that satisfies the conditions in \Cref{thm:wwp_dim_condition}. From these polytope conditions, it is quick to show that there is some finite sample threshold below which Fr\'echet means are almost surely not unique, and above which Fr\'echet means have a positive probability of uniqueness. This is proven by the following corollary.
	
	\begin{corollary}\label{cor:inductive_result}
		Suppose $U_{\CType}$ has positive probability, and that $F_{n+1}$ is a facet of $B$ such that $F_1 \subseteq F_{n+1}$. Then $U_{F_1,\dots,F_{n+1}}$ has positive probability. In particular, if $U_{\CType}$ has positive probability and induces a unique Fr\'echet mean, then $U_{F_1,\dots,F_{n+1}}$ has positive probability and induces a unique Fr\'echet mean.
	\end{corollary}
	
	\begin{proof}
		We first note that as $F_1 \subseteq F_{n+1}$, then $G_{n+1}$ must be some vertex of $G_1$, and so we have that $\conv \bigcup_{i \in [n+1]} G_i = \conv \bigcup_{i \in [n]} G_i$. Hence:
		\begin{align*}
			\vec 0 \in \relint \left( \conv \left(\textstyle \bigcup_{i \in [n]} G_i \right) \right) &\Rightarrow \vec 0 \in \relint \left( \conv \left(\textstyle \bigcup_{i \in [n+1]} G_i \right) \right), \\
			\dim \conv \left(\textstyle \bigcup_{i \in [n]} G_i \right) = k &\Rightarrow \dim \conv \left(\textstyle \bigcup_{i \in [n+1]} G_i \right) = k.
		\end{align*}
		Additionally, if $F_{i+1}$ is a facet, then $G_{i+1}$ and $\aff G_{i+1}$ are a single point. Therefore:
		\begin{align*}
			\textstyle \sum_{i \in [n+1]} \dim G_i &= 0 + \textstyle\sum_{i \in [n]} \dim G_i,\\
			&= \textstyle\dim \sum_{i \in [n]} \aff G_i,\\
			&= \textstyle\dim \sum_{i \in [n+1]} \aff G_i. 
		\end{align*}
		We conclude that if $U_{\CType}$ satisfies \cref{eq:wpp_cond_possible,eq:wpp_cond_probable,eq:wpp_cond_unique}, then so does $U_{F_1,\dots,F_n}$. Similarly, if $U_{\CType}$ induces a unique Fr\'echet mean, then so does $U_{F_1,\dots,F_n}$.
	\end{proof}
	
	As the probability of a unique Fr\'echet mean converges to one (\cref{thm:uniqueness_prob_convergence}), there is some finite sample size which gives a positive probability of a unique Fr\'echet mean. By \cref{cor:inductive_result}, we can then conclude that for any polytope norm, there is a finite sample threshold $N(B)$ such that a sample has a positive probability of producing a unique Fr\'echet mean if and only if $n\geq N(B)$. 
	
	\begin{definition}[Unique Fr\'echet mean sample threshold]
		Fix a polytope norm on $\R^k$ with ball $B$. The \emph{unique Fr\'echet mean sample threshold}, $N(B)$, is the smallest $n$ such that there is a positive probability of $\sample$ having a unique Fr\'echet mean.
	\end{definition}
	
	In the following section, we prove an upper bound for the unique Fr\'echet mean sample threshold $N(B)$ of a general polytope norm, and compute it explicitly for the $\ell_\infty$ and $\ell_1$ norms.
	
	\section{Calculating Uniqueness Sample Thresholds} \label{sec:calculating_examples}
	
	In this section we show the theoretical power of \cref{thm:wwp_dim_condition} by proving that for any polytope norm on $\R^k$, the unique Fr\'echet mean sample threshold is at most $k+1$. We then compute the unique Fr\'echet mean sample thresholds of the $\ell_\infty$ and $\ell_1$ norms --- which are $k$ and 3 respectively.
	
	\begin{proposition}\label{prop:threshold_bound}
		Consider a polytope normed space $(\R^k, \|\cdot\|_B)$ with unit ball $B$. The unique Fr\'echet mean sample threshold $N(B)$ is at most $k+1$.
	\end{proposition}
	
	\begin{proof}
		It suffices to find some $G_1,\dots, G_{k+1}$, each a face of $B^{\Delta}$, which together satisfy the conditions of \cref{thm:wwp_dim_condition}. We will construct $G_1,\dots, G_{k+1}$ iteratively, using the following result:
		
		\begin{lemma}[Corollary 11.7 of Chapter 2,  \cite{brondsted2012introduction}]\label{lem:vertex_neighbours}
			Let $Q \subset \R^k$ be some $k$-dimensional polytope with vertex $\vec x_0$, and let $\vec x_1, \dots, \vec x_r$ be the vertices adjacent to $\vec x_0$. Then
			\[
			\aff \{ \vec x_0,\dots, \vec x_r \} = \R^k.
			\]
		\end{lemma}
		We first fix $\vec v_1$ to be any vertex of the polar polytope $B^{\Delta}$. Let $G_1$ be $\vec v_1$, then we define vertices $\vec v_2, \dots, \vec v_k$ and faces $G_2, \dots, G_{k}$ iteratively like so:
		\begin{enumerate}
			\item Using the final line of \cref{lem:vertex_neighbours}, we fix $\vec v_{i+1}$ to be some neighbour of $-\vec v_i$ which is not in $\text{span} \{\vec v_1,\dots, \vec v_i\}$.
			\item Let $G_{i+1}$ be the edge connecting $(-\vec v_i, \vec v_{i+1})$.
		\end{enumerate}
		Note that by construction, $\vec v_1,\dots \vec v_k$ are linearly independent, so $\text{span}(\vec v_1, \dots, \vec v_{k-1})$ is some hyperplane not containing $\vec v_k$. There is then some vertex $\vec v_{k+1}$ of $B^{\Delta}$ which is on the other side of the hyperplane $\text{span}(\vec v_1, \dots, \vec v_{k-1})$ to $\vec v_k$. Let $G_{k+1}$ be the vertex $\vec v_{k+1}$.
		
		Having constructed $G_1, \dots, G_{k+1}$, it remains to prove that they satisfy \cref{eq:wpp_cond_possible,eq:wpp_cond_probable,eq:wpp_cond_unique}. We can first prove \cref{eq:wpp_cond_unique} directly, as $\vec v_1, \dots \vec v_k$ are linearly independent:
		\begin{align*}
			\dim \conv \left( \textstyle\bigcup_{i \in [k+1]} G_i  \right) &= \dim \aff \{ \pm \vec v_1, \dots, \pm \vec v_{k-1}, \vec v_k, \vec v_{k+1} \}, \\
			&= \dim \aff \{\vec 0, \vec v_0, \vec v_1, \dots, \vec v_k, \vec v_{k+1} \}, \\
			&= k.
		\end{align*}
		Next, we prove \cref{eq:wpp_cond_probable}. The edges $G_2, \dots, G_k$ each have dimension $1$, while $G_1, G_{k+1}$ are single vertices, so $\sum_{i \in [k+1]} \dim G_i = k-1$. Then:
		\begin{align*}
			\dim \left(\textstyle\sum_{i \in [k+1]} \aff G_i \right) &= \dim \left(\vec v_{k+1} + \vec v_1 + \textstyle\sum_{i \in [k-1]} (-\vec v_i + \text{span}(\vec v_i + \vec v_{i+1})) \right),\\
			&= \dim \text{span}(\vec v_1 + \vec v_2, \dots, \vec v_{k-1} + \vec v_{k}),\\
			&= k-1 = \textstyle \sum_{i \in [k+1]} \dim G_i,
		\end{align*}
		where the final line follows the fact that $\vec v_1 + \vec v_2, \dots, \vec v_{k-1} + \vec v_{k}$ must be independent as $\vec v_1, \dots, \vec v_k$ are independent. We have proved \cref{eq:wpp_cond_probable}. 
		
		It remains to prove \cref{eq:wpp_cond_possible}. It suffices to show that $\vec 0$ can be expressed as a positive sum of all vertices $\pm \vec v_1, \dots, \pm \vec v_{k-1}, \vec v_k, \vec v_{k+1}$. By definition, the vertices $\vec v_k$ and $\vec v_{k+1}$ are on opposite sides of the hyperplane $\text{span}(\vec v_1, \dots, \vec v_{k-1})$, so there are some $\lambda_1, \dots, \lambda_{k-1}$ and positive $\mu_k, \mu_{k+1}$ such that 
		\begin{align*}
			\vec 0 &= \mu_k \vec v_k + \mu_{k+1} \vec v_{k+1} + \textstyle\sum_{i \in [k-1]} \lambda_i \vec v_i, \\
			&= \mu_k \vec v_k + \mu_{k+1} \vec v_{k+1} + \textstyle\sum_{i \in [k-1]} (\lambda_i + |\lambda_i|+1) \vec v_i + (|\lambda_i| + 1)(-\vec v_i).
		\end{align*}
		This proves \cref{eq:wpp_cond_possible}, and therefore $G_1, \dots, G_{k+1}$ satisfy \cref{eq:wpp_cond_possible,eq:wpp_cond_probable,eq:wpp_cond_unique}. The unique Fr\'echet mean sample threshold $N(B)$ is therefore at most $k+1$.
	\end{proof}
	
	\subsection{Uniqueness sample thresholds for the $\ell_{\infty}$ and $\ell_1$ norms} \label{subsec:UST_examples}
	
	Using the conditions of \Cref{thm:wwp_dim_condition}, we can identify the uniqueness sample threshold in the case of the $\ell_{\infty}$ and $\ell_1$ norms.  \\
	
	We first show that the $\ell_\infty$ norm has a unique Fr\'echet mean sample threshold of $k$ for $k\geq 3$, and a unique Fr\'echet mean sample threshold of $3$ for $k= 2$. The face construction used in the proof is show for the $k=3$ case in \cref{fig:hypercube_face_types}.
	
	\begin{proposition}\label{prop:hypercube}
		Consider the normed space $(\R^k, \ell_{\infty})$, so $B=[-1,1]^k$. When $k \geq 3$, the sample threshold $N([-1,1]^k)$ is $k$. In the $k = 2$ case, $N([-1,1]^2) = 3$.
	\end{proposition}
	
	\begin{proof}
		We first show that if $n<k$, then $G_1, \dots, G_n$ cannot satisfy \cref{eq:wpp_cond_possible,eq:wpp_cond_probable,eq:wpp_cond_unique}.
		
		Suppose $G_1, \dots, G_n$ satisfy \cref{eq:wpp_cond_possible,eq:wpp_cond_probable,eq:wpp_cond_unique}. As $B = [-1, 1]^k$, the polar polytope $B^{\Delta}$ is given by the cross-polytope $\conv \{ \pm \vec e_1,\dots, \pm \vec e_k \}$. Let $\mathscr V$ be the vertex set of $B^{\Delta}$; that is, $\mathcal V = \{ \pm \vec e_1,\dots, \pm \vec e_k \}$.
		
		If any one of the unit vectors $\vec e_j$ is not in $\bigcup_{i \in [n]} G_i$, then $\conv \left( \bigcup_{i \in [n]} G_i \right)$ is in the half-space with non-positive $j^{th}$ coordinate. As $k = \dim \conv \left( \bigcup_{i \in [n]} G_i\right)$ by \cref{eq:wpp_cond_unique}, we then conclude $\vec 0 \notin \relint \left(\conv \left( \bigcup_{i \in [n]} G_i\right)\right)$; contradiction. Similarly, $\bigcup_{i \in [n]} G_i$ must also contain every $-\vec e_j$. Hence $\bigcup_{i \in [n]} G_i$ contains every vertex in $\mathscr V$.
		
		Every face of $B^{\Delta}$ is a simplex, so a face $G_i$ of dimension $m$ contains $m+1$ vertices of $B^{\Delta}$. By \cref{eq:wpp_cond_probable}:
		\begin{align*}
			k &\geq \dim \textstyle\sum_{i \in [n]} \aff G_i, \\
			&= \textstyle \sum_{i \in [n]} \dim G_i, \\
			&= \textstyle \sum_{i \in [n]} \left(|G_i \cap \mathscr V| - 1\right), \\
			&\geq \textstyle 2k - n.
		\end{align*}
		Hence $n \geq k$; contradiction. 
		
		We now show that for $k \geq 3$, there exists $G_1, \dots, G_k$ satisfying \cref{eq:wpp_cond_possible,eq:wpp_cond_probable,eq:wpp_cond_unique}. We define 
		$$
		G_1, \dots, G_k = \conv (\vec e_1,-\vec e_2), \dots, \conv (\vec e_k, -\vec e_1).
		$$
		This construction is shown in \cref{fig:hypercube_face_types} for $k=3$.
		
		Then $\conv \bigcup_{i \in [k]} G_i = B^{\Delta}$, which has dimension $k$ and $\vec 0$ in its interior. Also:
		\begin{align*}
			\dim \textstyle \sum_{i \in [k]} \aff G_i &= \dim \text{span} \{ e_1 + e_2, \dots, e_k + e_1 \}, \\
			&= k = \textstyle \sum_{i \in [k]} \dim G_i.
		\end{align*}
		In computing the span dimension, we used $k \geq 3$. We conclude $N([-1,1]^k) = k$ for $k \geq 3$.
		
		We now consider the $k=2$ case. By \cref{rmk:non_unique_for_2n}, $N([-1,1]^2) > 2$. We can verify that $G_1 = -\vec e_1, G_2 = -\vec e_2, G_3 = \conv \{\vec e_1, \vec e_2 \}$ satisfies \cref{eq:wpp_cond_possible,eq:wpp_cond_probable,eq:wpp_cond_unique}, so $N([-1,1]^2) = 3$.    
	\end{proof}
	
	\begin{figure*}[ht!]
		\centering
		\begin{subfigure}[b]{0.48\textwidth}
			\centering
			\tdplotsetmaincoords{70}{50}
			\begin{tikzpicture}[scale=1.33, tdplot_main_coords]
				\coordinate (O) at (0,0,0);
				
				\draw[->] (-3,0,0) -- (3,0,0) node[anchor=north east]{$x$};
				\draw[->] (0,3,0) -- (0,-3,0) node[anchor=north west]{$y$};
				\draw[->] (0,0,-2.25) -- (0,0,2.25) node[anchor=south]{$z$};
				
				\coordinate (A) at (1,1,1);
				\coordinate (B) at (1,-1,1);
				\coordinate (C) at (-1,-1,1);
				\coordinate (D) at (-1,1,1);
				\coordinate (E) at (1,1,-1);
				\coordinate (F) at (1,-1,-1);
				\coordinate (G) at (-1,-1,-1);
				\coordinate (H) at (-1,1,-1);
				
				\draw (A) -- (B) -- (C) -- (D) -- (A);
				\draw (E) -- (F) -- (G);
				\draw[dashed] (G) -- (H);
				\draw[dashed] (E) -- (H);
				\draw (A) -- (E);
				\draw (B) -- (F);
				\draw (C) -- (G);
				\draw[dashed] (D) -- (H);
				\fill[fill=gray,opacity=1/8](A) -- (B) -- (F) -- (E) -- cycle;
				\fill[fill=gray,opacity=1/8](C) -- (D) -- (H) -- (G) -- cycle;
				\fill[fill=gray,opacity=1/8](A) -- (D) -- (H) -- (E) -- cycle;
				\fill[fill=gray,opacity=1/8](C) -- (B) -- (F) -- (G) -- cycle;
				\fill[fill=gray,opacity=1/8](A) -- (B) -- (C) -- (D) -- cycle;
				\fill[fill=gray,opacity=1/8](E) -- (F) -- (G) -- (H) -- cycle;
				
				\draw[red,line width=3,line cap=round,opacity=3/4] (A) -- (E)  node[midway,anchor=west]{$F_1$};
				\draw[red,line width=3,line cap=round,opacity=3/4] (D) -- (C)  node[midway,anchor=south]{$F_3$};
				\draw[red,line width=3,line cap=round,opacity=3/4] (G) -- (F)  node[midway,anchor=north]{$F_2$};
				
			\end{tikzpicture}
			\caption{The 3-dimensional hypercube, which is the unit ball for the $\ell_{\infty}$ norm, and a face type which induces a unique Fr\'echet mean.}
		\end{subfigure}
		\hfill
		\begin{subfigure}[b]{0.48\textwidth}
			\centering
			\tdplotsetmaincoords{70}{50}
			\begin{tikzpicture}[scale = 2,tdplot_main_coords]
				\coordinate (O) at (0,0,0);
				
				\draw[->] (-2,0,0) -- (2,0,0) node[anchor=north east]{$x$};
				\draw[->] (0,2,0) -- (0,-2,0) node[anchor=north west]{$y$};
				\draw[->] (0,0,-1.5) -- (0,0,1.5) node[anchor=south]{$z$};
				
				\coordinate (A) at (1,0,0);
				\coordinate (B) at (0,1,0);
				\coordinate (C) at (0,0,1);
				\coordinate (D) at (-1,0,0);
				\coordinate (E) at (0,-1,0);
				\coordinate (F) at (0,0,-1);
				
				\draw (D) -- (E) -- (A) -- (B);
				\draw[dashed] (D) -- (B);
				\draw (D) -- (C) -- (A) -- (F);
				\draw[dashed] (D) -- (F);
				\draw (F) -- (E) -- (C) -- (B);
				\draw[dashed] (F) -- (B);
				\fill[fill=gray,opacity=1/8](A) -- (B) -- (C) -- cycle;
				\fill[fill=gray,opacity=1/8](D) -- (E) -- (F) -- cycle;
				\fill[fill=gray,opacity=1/8](A) -- (B) -- (F) -- cycle;
				\fill[fill=gray,opacity=1/8](C) -- (D) -- (E) -- cycle;
				\fill[fill=gray,opacity=1/8](A) -- (E) -- (C) -- cycle;
				\fill[fill=gray,opacity=1/8](D) -- (B) -- (F) -- cycle;
				\fill[fill=gray,opacity=1/8](A) -- (E) -- (F) -- cycle;
				\fill[fill=gray,opacity=1/8](D) -- (B) -- (C) -- cycle;
				
				\draw[red,line width=3,line cap=round,opacity=3/4] (B) -- (A)  node[midway,anchor=west]{$G_1$};
				\draw[red,line width=3,line cap=round,opacity=3/4] (D) -- (C)  node[midway,anchor=base east]{$G_3$};
				\draw[red,line width=3,line cap=round,opacity=3/4] (F) -- (E)  node[midway,anchor=east]{$G_2$};
			\end{tikzpicture}
			\caption{The polar polytope for the 3-dimensional hypercube, and the faces polar to $F_1, F_2, F_3$ which satisfy conditions \cref{eq:wpp_cond_possible,eq:wpp_cond_probable,eq:wpp_cond_unique} in \cref{thm:wwp_dim_condition}.}
		\end{subfigure}
		\caption{A face type for the $\ell_{\infty}$ unit ball in 3 dimensions which induces a unique Fr\'echet mean. In the proof of \cref{prop:hypercube}, we use a generalisation of this example to show the unique Fr\'echet mean sample threshold for the $\ell_{\infty}$ norm in $\R^k$ is $k$ for $k\geq 3$.}
		\label{fig:hypercube_face_types}
	\end{figure*}
	
	When using the $\ell_{\infty}$ norm, the sample size threshold is the dimension of the space, $k$. We might expect this to be standard behaviour, but in fact for the $\ell_1$ norm the sample size threshold is just 3.
	
	\begin{proposition}\label{prop:cross-polytope}
		Consider the normed space $(\R^k, \ell_1)$, so $B=\conv \{ \pm \vec e_1, \dots, \pm \vec e_k \}$. Then $N(B) = 3$.
	\end{proposition}
	
	\begin{proof}
		The polar polytope of the $\ell_1$ ball is the hypercube. We fix $G_1$ to be the face polar to $-\vec e_1$, so $G_1 = \{ (-1, x_2, \dots, x_k): |x_i| \leq 1 \}$. Now let $G_2$ be the vertex $(1,\dots, 1)$ and let $G_3$ be the edge between $ (1, -1, \dots, -1)$ and $(-1, -1, \dots, -1)$. 
		
		The face $G_1$ has dimension $k-1$ but $G_2$ is not in the affine hull of $G_1$, so their convex hull must have dimension $k$. This proves \cref{eq:wpp_cond_unique}.
		
		The affine hulls of $G_1, G_2, G_3$ are $\{ \vec x : x_1 = -1 \}, \{ (1,\dots, 1) \}, \{ (\lambda, -1, \dots, -1): \lambda \in \R \}$ respectively. The sum of their affine hulls is the entire space $\R^k$ and so has dimension $k$, which is also the sum of their individual dimensions. This proves \cref{eq:wpp_cond_probable}. 
		
		Note that $(-1, 0,\dots, 0)$ is in the relative interior of $G_1$, $(1,\dots, 1)$ is the relative interior of $G_2$, and $(0, -1, \dots, -1)$ is in the relative interior of $G_3$. We then have:
		\[
		\vec 0 = \tfrac{1}{3} (-1, 0,\dots, 0) + \tfrac{1}{3} (1,\dots, 1) + \tfrac{1}{3} (0, -1, \dots, -1)
		\] 
		Hence $\vec 0 \in \relint \conv (G_1 \cup G_2 \cup G_3)$, and we have proved \cref{eq:wpp_cond_possible}.
		
		By \cref{rmk:non_unique_for_2n}, we have $N(B)>2$, so $N(B) = 3$.
	\end{proof}
	
	\begin{remark}
		We note that the construction in \cref{prop:cross-polytope} can be replicated for any polytope whose polar $B^{\Delta}$ has some facet $G$ with vertices $\vec u$ and $\vec v$ whose mid-point passes through the relative interior of $G$. This includes any norm whose polar is a zonotope.
	\end{remark}
	
	\begin{figure*}[ht!]
		\centering
		\begin{subfigure}[b]{0.48\textwidth}
			\centering
			\tdplotsetmaincoords{70}{40}
			\begin{tikzpicture}[scale = 2,tdplot_main_coords]
				\coordinate (O) at (0,0,0);
				
				\draw[->] (-2,0,0) -- (2,0,0) node[anchor=north east]{$x$};
				\draw[->] (0,2,0) -- (0,-2,0) node[anchor=north west]{$y$};
				\draw[->] (0,0,-1.5) -- (0,0,1.5) node[anchor=south]{$z$};
				
				\coordinate (A) at (1,0,0);
				\coordinate (B) at (0,1,0);
				\coordinate (C) at (0,0,1);
				\coordinate (D) at (-1,0,0);
				\coordinate (E) at (0,-1,0);
				\coordinate (F) at (0,0,-1);
				\draw[red,line width=3,line cap=round,opacity=3/4] (B) -- (F)  node[anchor=north west]{$F_3$};
				
				\draw (D) -- (E) -- (A) -- (B);
				\draw[dashed] (D) -- (B);
				\draw (D) -- (C) -- (A) -- (F);
				\draw[dashed] (D) -- (F);
				\draw (F) -- (E) -- (C) -- (B);
				\draw[dashed] (F) -- (B);
				\fill[fill=gray,opacity=1/8](A) -- (B) -- (C) -- cycle;
				\fill[fill=gray,opacity=1/8](D) -- (E) -- (F) -- cycle;
				\fill[fill=gray,opacity=1/8](A) -- (B) -- (F) -- cycle;
				\fill[fill=gray,opacity=1/8](C) -- (D) -- (E) -- cycle;
				\fill[fill=gray,opacity=1/8](A) -- (E) -- (C) -- cycle;
				\fill[fill=gray,opacity=1/8](D) -- (B) -- (F) -- cycle;
				\fill[fill=gray,opacity=1/8](A) -- (E) -- (F) -- cycle;
				\fill[fill=gray,opacity=1/8](D) -- (B) -- (C) -- cycle;
				
				
				\filldraw[red,opacity=3/4] (D) circle (1pt) node[anchor=east]{$F_1$};
				\fill[fill=red,opacity=3/4](A) -- (E) -- (C) node[anchor=south west,red]{$F_2$} -- cycle;
			\end{tikzpicture}
			\caption{The 3-dimensional cross-polytope, which is the unit ball for the $\ell_{1}$ norm, and a face type which induces a unique Fr\'echet mean.}
		\end{subfigure}
		\hfill
		\begin{subfigure}[b]{0.48\textwidth}
			\centering
			\tdplotsetmaincoords{70}{40}
			\begin{tikzpicture}[scale=1.33, tdplot_main_coords]
				\coordinate (O) at (0,0,0);
				
				\draw[->] (-3,0,0) -- (3,0,0) node[anchor=north east]{$x$};
				\draw[->] (0,3,0) -- (0,-3,0) node[anchor=north west]{$y$};
				\draw[->] (0,0,-2.25) -- (0,0,2.25) node[anchor=south]{$z$};
				
				\coordinate (A) at (1,1,1);
				\coordinate (B) at (1,-1,1);
				\coordinate (C) at (-1,-1,1);
				\coordinate (D) at (-1,1,1);
				\coordinate (E) at (1,1,-1);
				\coordinate (F) at (1,-1,-1);
				\coordinate (G) at (-1,-1,-1);
				\coordinate (H) at (-1,1,-1);
				
				\draw (A) -- (B) -- (C) -- (D) -- (A);
				\draw (E) -- (F) -- (G);
				\draw[dashed] (G) -- (H);
				\draw[dashed] (E) -- (H);
				\draw (A) -- (E);
				\draw (B) -- (F);
				\draw (C) -- (G);
				\draw[dashed] (D) -- (H);
				\fill[fill=gray,opacity=1/8](A) -- (B) -- (F) -- (E) -- cycle;
				\fill[fill=gray,opacity=1/8](C) -- (D) -- (H) -- (G) -- cycle;
				\fill[fill=gray,opacity=1/8](A) -- (D) -- (H) -- (E) -- cycle;
				\fill[fill=gray,opacity=1/8](C) -- (B) -- (F) -- (G) -- cycle;
				\fill[fill=gray,opacity=1/8](A) -- (B) -- (C) -- (D) -- cycle;
				\fill[fill=gray,opacity=1/8](E) -- (F) -- (G) -- (H) -- cycle;
				
				\filldraw[color=red,fill=red,opacity=3/4] (B) circle (1pt) node[anchor=west,red]{$G_2$};
				\fill[fill=red,opacity=3/4](C) node[anchor=south east,red]{$G_1$} -- (D) -- (H) -- (G) -- cycle;
				\draw[red,line width=3,line cap=round,opacity=3/4] (E) -- (H)  node[midway,anchor=north]{$G_3$};
			\end{tikzpicture}
			\caption{The polar polytope for the 3-dimensional cross-polytope, and the faces polar to $F_1, F_2, F_3$ which satisfy conditions \cref{eq:wpp_cond_possible,eq:wpp_cond_probable,eq:wpp_cond_unique} in \cref{thm:wwp_dim_condition}.}
		\end{subfigure}
		\caption{A face type for the $\ell_{1}$ unit ball in 3 dimensions which induces a unique Fr\'echet mean. In the proof of \cref{prop:cross-polytope}, we use a generalisation of this example to show the unique Fr\'echet mean sample threshold for the $\ell_{1}$ norm in $\R^k$ is $3$.}
		\label{fig:cross-polytope_face_types}
	\end{figure*}
	
	\cref{fig:cross-polytope_face_types} shows our construction of a 3-sample face type which has a positive probability of producing a unique Fr\'echet mean in the case of the $\ell_1$ norm in 3 dimensions. \\
	
	In this section we have used \cref{thm:wwp_dim_condition} to study the unique Fr\'echet mean sample threshold, not only computing the sample thresholds for the $\ell_\infty$ and $\ell_1$ norms, but also proving a general upper bound of $k+1$, where $k$ is the dimension of our state space.
	
	\section{Exact Computation of Fr\'echet Means} \label{sec:computation}
	
	We conclude this paper with a brief note on computation. In order to run numerical experiments on the dimension of a Fr\'echet set, we must be able to perform exact computations.
	
	We have defined Fr\'echet means as the solutions of the global minimisation of a non-smooth objective function, but we can re-frame this as the constrained optimisation of a quadratic form.
	
	\begin{lemma}\label{lem:constrained_quadratic_optimisation}
		A point $\theta$ is in $\FMset(\vec x_1,\dots, \vec x_n)$ if and only if there is some $(d_1, \dots, d_n)$ such that $\theta, \vec d$ give a solution to the following constrained quadratic optimisation problem: 
		\begin{align} \label{eq:FM_objective}
			\text{minimise}\quad d_1^2+\dots+d_n^2, \qquad \text{given} \quad \forall \;  i \in [n], \vec a \in A: \quad \vec a^{\T}\vec x_i \leq d_i +\vec a^{\T}\theta.
		\end{align}    
	\end{lemma}
	\begin{proof}
		Let $Q$ be the polyhedron defined by $Q = \{(\theta, \vec d) : \forall i \in [n], \vec a \in A , \, \vec a^\T(\vec x_i - \theta) \leq d_i \}$. Notice that any $(\theta, \vec d)$ in $Q$ which minimises \cref{eq:FM_objective} will be of the form $(\theta, (d(\theta,\vec x_i))_{i \in [n]})$, as otherwise there is some $d_i$ we can decrease. Hence minimisation of $\|\vec d\|_2^2$ over $Q$ is equivalent to
		\begin{align*}
			\text{minimise}&\quad d_1^2+\dots+d_n^2\\
			\text{given} &\quad \forall \;  i \in [n], d_i = \|\vec x_i - \theta\|_B.
		\end{align*}
		This is exactly the optimisation problem of sample Fr\'echet means.
	\end{proof}
	We have translated the tropical Fr\'echet mean problem into a quadratic optimisation problem with linear constraints in $\R^{n+k}$. In order to compute the full Fr\'echet mean set, and in particular to study its dimension, we note that \cref{lem:constrained_quadratic_optimisation} allows us to write out $\FMset$ more explicitly. Let $Q \subset \R^{k+n}$ be the polyhedron given by $Q = \{(\theta, \vec d) : \forall i \in [n], \vec a \in A , \, \vec a^\T(\vec x_i - \theta) \leq d_i \}$, and let $Q'$ be the projection of $Q$ to $\R^n$. Then the Fr\'echet mean set $\FMset(\sample)$ is given by: 
	\[
	\FMset = \pi_{\R^k}\pi^{-1}_{\R^n}(\argmin_{\vec d \in Q'} \norm{2}{\vec d})
	\]
	
	This observation enables us to break down our computation of $\FMset$ into three steps:
	\begin{enumerate}
		\item Compute a representation of $Q'$,
		\item Compute the optimal $\vec d' = \argmin_{\vec d \in Q'} \norm{2}{\vec d}$,
		\item Compute the polytope $\FMset = \pi_{\R^k}\pi^{-1}_{\R^n} \vec d'$ .
	\end{enumerate}
	
	Step 3 is routine; from $\vec d'$, we can write down the hyperplane representation for $\FMset(\sample)$ as
	\[
	\FMset(\sample) = \{ \theta \in \R^k : \, \forall \vec a \in A, \, \vec a^\T \vec x_i - d_i' \leq \vec a^\T \theta \}.
	\]
	Step 2 is the most computationally intensive part of this computation, but can be solved exactly using the Frank-Wolfe \citep{frank1956algorithm} algorithm. This algorithm applies iterated linear programming to find a point in a polytope in the positive quadrant with minimal Euclidean norm; this is exactly the problem of computing $\vec d' = \argmin_{\vec d \in Q} \norm{2}{\vec d}$. The Frank-Wolfe algorithm avoids having to compute matrix inverses and is known to terminate in $n$ iterations, where $n$ is the dimension of the ambient space.
	
	Step 1 requires projecting $Q \subset \R^{n+k}$ to $Q' \subset\R^n$. The initial polyhedron $Q$ is given in hyperplane form, so to project this to $\R^n$, we perform Fourier--Motzkin elimination \cite[Chapter 1]{ziegler2012lectures} to compute the hyperplane representation of $Q'$ from the hyperplane representation of $Q$.
	
	This full procedure is presented in \cref{alg:exact_computation}.
	
	\begin{algorithm}
		\caption{Fr\'echet Mean Set Computation}\label{alg:exact_computation}
		\begin{algorithmic}
			\Require normal vector matrix $A$, data matrix $\vec x$.
			\State $Q \gets \{(\theta, \vec d) : \forall i \in [n], \vec a \in A , \, \vec a^\T \vec x_i - d_i \leq \vec a^\T \theta \}$\Comment{Feasible region for $(\theta, \vec d)$ in $\mathbb{R}^{k+n}$}
			\State $Q' \gets \text{Fourier--Motzkin}(Q, [k+1:k+n])$ \Comment{Feasible region for $\vec d$ in $\mathbb{R}^n$}
			\State $d_0 \gets \argmin_{d \in Q'} \, \vec 1^\T \vec d$
			\State $t=0$
			\While{$t \leq n$} \Comment{Perform $n$ steps of Frank-Wolfe}
			\State $\vec s_{t} \gets \argmin_{\vec s \in Q'} \vec s^{\T}\vec d_t$
			\State $\eta_t \gets \min \{-(\vec s_t - \vec d_t)^{\T} \vec d_t / 2\|\vec s_t - \vec d_t\|_2^2 , 1 \}$
			\State $\vec d_{t+1} \gets \vec d_t + \eta_t(\vec s_t-\vec d_t)$
			\State $t \gets t+1$
			\EndWhile
			\State $\vec d' \gets \vec d_n$
			\State $\FMset \gets \{ \theta \in \R^k : \, \forall \vec a \in A, \, \vec a^\T \vec x_i - d_i' \leq \vec a^\T \theta \}$ \Comment{Compute Fr\'echet mean set}\\
			\Return $\FMset$
		\end{algorithmic}
	\end{algorithm}
	
	We now use \cref{alg:exact_computation} to verify our results on the sample threshold for unique Fr\'echet means by showing that the sample threshold for uniqueness is dependent on the dimension $k$ for $\ell_{\infty}$ Fr\'echet means.
	
	\begin{figure}[ht]
		\centering
		\includegraphics[width=0.6\linewidth]{Linf_uniqueness_probability.png}
		\caption{100 independent datasets of rationalised Gaussian data of size $n$ in dimension $k$ are generated. This shows the proportion of the 100 datasets which have a unique $\ell_{\infty}$-Fr\'echet mean.}
		\label{fig:Linf_uniqueness_probability}
	\end{figure}
	
	When computing $\ell_{\infty}$ Fr\'echet means, we run dimension computations for $2 \leq n,k \leq 10$; this is less computationally intensive compared to other norms, in part due to the small number of normal vectors and their linear independence. For each pair $k,n$, we generate 100 independent sets of rationalised Gaussian data and compute the dimension of each Fr\'echet mean set. \Cref{fig:Linf_uniqueness_probability} shows the proportion of these datasets producing a unique Fr\'echet mean as the sample size $n$ increases.
	
	In \cref{fig:Linf_uniqueness_probability}, we see that the uniqueness probability is initially 0 before roughly increasing. The value of $n$ at which this probability becomes positive is increasing in $k$; for $k=2,3$, the observed sample threshold is $3$, as we proved in \cref{prop:hypercube}. However, for higher $k$ the observed sample threshold is slightly higher than expected. This is likely due to the fact that even though $p_k$ is theoretically positive, it may be a very low possibility so not observable over 100 samples.
	
	\section{Discussion}
	
	In this paper we have shown that the probability of a unique Fr\'echet mean over a polytope normed space converges to 1, and we have identified geometric conditions which govern the number of sample points needed before Fr\'echet means are unique with positive probability. In doing so, we were able to define the \emph{unique Fr\'echet mean sample threshold} of a polytope norm --- Fr\'echet means are almost never unique below this sample threshold, but above this threshold the probability of uniqueness is positive and converging to one. We have shown this threshold to be at most $k+1$, where $k$ is the dimension of our space, and computed it exactly for the $\ell_\infty$ and $\ell_1$ norms. Finally, we presented an algorithm for exact computation, which can be used to study the probability of unique Fr\'echet means through simulation.
	
	We have shown an upper bound of $k+1$ for the unique Fr\'echet mean sample threshold of a general polytope norm, which is within one of our $\ell_{\infty}$ example. We hope to close this gap as future work, finding a strict upper bound. This will require either finding a polytope example with a unique Fr\'echet mean sample threshold of $k+1$, or lowering our general upper bound to $k$ and showing the $\ell_\infty$ case to be the worst possible.
	
	Currently, our geometric conditions which dictate the sample threshold for unique Fr\'echet means (\cref{thm:wwp_dim_condition}) do not provide a way of computing the threshold for an arbitrary polytope. At best, one could test every $n$-tuple of faces in $\mathfrak F (B^{\Delta})$ against the conditions in \cref{thm:wwp_dim_condition} for increasing $n$ until the necessary conditions are satisfied. It would be of interest to find a more direct method of computation, perhaps utilising the Gale transform of $B^{\Delta}$ \cite[Chapter 6]{ziegler2012lectures}, which records which sets of vertices contain $\vec 0$ in the relative interior of their convex hull. Relating the sample threshold of $B$ to known values in matroid theory would be particularly useful for finding efficient methods of computation for sample thresholds.
	
	The conditions given in \cref{thm:wwp_dim_condition} state which face types have a positive probability of occurring, and independently, which face types give rise to a unique Fr\'echet mean. We can extend this to get conditions for face types to produce a $d$-dimensional Fr\'echet mean set with positive probability, and then we can adapt \cref{cor:inductive_result} to formalise a $d$-dimensional Fr\'echet mean sample threshold. This raises the question - how do the sample thresholds compare for different values of $d$? As an initial investigation, we can study their behaviour computationally using \cref{alg:exact_computation}.
	
	\section{Acknowledgements}
	A.Mc. acknowledges the support of the Natural Sciences and Engineering Research Council of Canada (NSERC), [RGPIN-2025-03968, DGECR-2025-00237]. A.Mo. and R.T. are partly funded by EPSRC grant number EP/Y028872/1 Mathematical Foundations of Intelligence.
	
	\newpage
	
	\bibliographystyle{authordate3}
	
	\bibliography{Sources}
	
	\pagebreak
	
	\appendix
	
	\section{Zero Subgradient Event Proof} \label{appsec:prop_proof}
	
	\EventEquivalence*
	
	The proof of Proposition \ref{prop:lifted_type_event} follows directly from the lemmas below.
	
	\begin{lemma}\label{lemma:polyhedral_cone_W}
		The set $W = \{\vec w \in \R^n_{+} : \vec 0 \in \sum_{i \in [n]} w_i \relint G_i\}$ is a relatively open, polyhedral cone.
	\end{lemma}
	
	\begin{proof}
		Let $\vec g_1^{(i)}, \dots, \vec g_{r_i}^{(i)} \in \R^k$ be the vertices of $G_i$. Let $M_G \in \R^{k \times \sum_{i \in [n]} r_i}$ be the matrix whose columns are given by all of the $\vec g_j^{(i)}, \, i\in [n], j \in [r_i]$. We then index $\vec x \in  \R^{\sum_{i \in [n]} r_i}$ by $\vec x = (x_1^{(1)},\ldots,x_{m_1}^{(1)},\ldots,x_{1}^{(n)},\ldots,x^{(n)}_{m_n})$. Define the linear map $L_W: \R^{\sum_{i \in [n]} r_i} \rightarrow \R^n$ by $\left(L_W(\vec x)\right)_i = \sum_{j \in [r_i]} x_j^{(i)}$. We will show that $W = L_W(\Ker M_G \cap \R^{\sum_{i \in [n]} r_i}_{+})$. \\
		By \cite[Theorem 6.9]{rockafellar1970convex}, which we quoted in \cref{eq:rockafellar_thm6.9}:
		\[
		\relint G_i = \bigg\{ \textstyle \sum_{j \in [r_i]} \lambda^{(i)}_{j} \vec g_j^{(i)} : \sum_{j \in [r_i]} \lambda^{(i)}_j = 1,\; \lambda^{(i)}_j > 0 \bigg\}.
		\]
		Therefore $\vec w \in W$ if and only if there exist positive $\left\{\lambda^{(i)}_j\right\}_{i \in [n], j \in [r_i]}$ satisfying $\sum_{j \in [r_i]} \lambda_j^{(i)} = 1$ and 
		\[
		\vec 0 = \textstyle \sum_{i \in [n]} \sum_{j \in [r_i]} w_i \lambda^{(i)}_j \vec g_j^{(i)}.
		\]
		For any $\vec w \in W$, pick $\vec \lambda$ satisfying the above and define $\vec x$ by $x_j^{(i)} = w_i\lambda_j^{(i)} > 0$. Then we have that $M_G \vec x = \vec 0$ and $L_W(\vec x) = \vec w$, and we conclude that $W \subseteq L_W(\Ker M_G \cap \R^{\sum_{i \in [n]} r_i}_{+})$. \\
		Now suppose $\vec x \in \Ker M_G \cap \R^{\sum_{i \in [n]} r_i}_{+}$. Then setting $\lambda_j^{(i)} = x_j^{(i)}\left(\sum_{j \in [r_i]} x_j^{(i)}\right)^{-1}$, we see that $L_W(\vec x)$ satisfies the condition above and so is in $W$. We conclude that $L_W(\Ker M_G \cap \R^{\sum_{i \in [n]} r_i}_{+}) \subseteq W$. \\
		As $\Ker M_G \cap \R^{\sum_{i \in [n]} r_i}_{+}$ is a relatively open, polyhedral cone, so is $L_W(\Ker M_G \cap \R^{\sum_i r_i}_{>0}) = W$.
	\end{proof}
	
	\begin{lemma}\label{lemma:RelintGrad_to_RelintFMset}
		Given any points  $(\sample, \theta) \in \R^{k(n+1)}$, if 
		$\vec 0 \in \relint(\partial f(\theta))$ then $\theta \in \relint \FMset$.
	\end{lemma}
	
	\begin{proof}
		We prove this lemma by contradiction; suppose that $\vec 0 \in \relint\partial f(\theta)$ but $\theta \notin \relint \FMset$. We begin by translating the property $\vec 0 \in \relint \partial f(\theta)$ into a condition on the directional derivatives of $f$.
		\begin{claim*}
			For all $\vec v \in \R^k$, either $\partial_{\vec v} f(\theta) = \partial_{- \vec v} f(\theta) = 0$ or $\partial_{\vec v} f(\theta), \partial_{-\vec v} f(\theta) > 0$.
		\end{claim*}
		\begin{proof}
			For all $\vec v \in (\partial f(\theta))^{\perp}$, we immediately have
			\begin{align*}
				\partial_{\vec v} f(\theta) = \sup_{\vec g \in \partial f(\theta)} \vec g^{\T}\vec v = 0, \qquad 
				\text{ and } \qquad \partial_{-\vec v} f(\theta) = \sup_{\vec g \in \partial f(\theta)} -\vec g^{\T}\vec v = 0.
			\end{align*}
			Now consider $\vec v \notin (\partial f(\theta))^{\perp}$, so there is some $\vec g_1 \in \partial f(\theta)$ such that $\vec g_1^{\T}\vec v \neq 0$. As $\vec 0 \in \relint\partial f(\theta)$, there is some $\vec g_2 \in \partial f(\theta)$ such that $\vec 0 = \lambda \vec g_1 + (1-\lambda) \vec g_2$, $0 < \lambda < 1$. Therefore $0=\lambda \vec g_1^{\T}\vec v + (1-\lambda) \vec g_2^{\T} \vec v$, and $\vec g_1^{\T}\vec v$ and $\vec g_2^{\T}\vec v$ must be non-zero with opposite signs. Hence:
			\begin{align*}
				\partial_{\vec v} f(\theta) = \sup_{\vec g \in \partial f(\theta)} \vec g^{\T}\vec v > 0, \qquad
				\text{ and } \qquad \partial_{-\vec v} f(\theta) = \sup_{\vec g \in \partial f(\theta)} -\vec g^{\T}\vec v > 0.
			\end{align*}
		\end{proof}    
		We have assumed that $\vec 0 \in \relint\partial f(\theta)$ but $\theta \notin \relint \FMset$, so $\theta$ is on the relative boundary of $\FMset$. Therefore there is some direction $\vec v$ such that $\theta + \lambda \vec v \in \FMset$ and $\theta - \lambda\vec v \notin \FMset$ for all positive $\lambda$. Then by the convexity of $f$ we have that $\partial_{\vec v} f(x) = 0$, and so by the claim above we also have $\partial_{-\vec v} f(x) = 0$.
		
		Note that $f$ is a maximum over quadratic functions whose coefficients are given by vectors $\vec a \in A$:
		\[
		f(\vec x) = \textstyle \frac{1}{n} \sum_{i \in [n]} \|X_i - \vec x \|_B^2 =  \frac{1}{n} \sum_{i \in [n]} \max_{\vec a \in A} \left[ \vec a^{\T}(X_i- \vec x) \right]^2 =  \max_{\vec a_1,\dots, \vec a_n \in A}  \frac{1}{n} \sum_{i \in [n]} \left[\vec a_i^{\T}(X_i- \vec x)\right]^2.
		\]
		Then we can fix $\vec a_1,\dots, \vec a_n, \vec b_1, \dots, \vec b_n$ to be the coefficients of quadratics $f_1, f_2$ which agree with $f$ on the intervals $[\theta, \theta+\epsilon \vec v]$ and $[\theta - \epsilon \vec v, \theta]$ respectively. Explicitly, we pick $\vec a_1,\dots, \vec a_n, \vec b_1, \dots, \vec b_n \in A$ such that for some positive epsilon:
		\begin{align*}
			\text{ for }\lambda \in [0,\epsilon]:  \qquad f(\theta + \lambda \vec v) = f_1(\lambda) &\coloneqq \textstyle \sum_{i \in [n]} \vec a_i^{\T}(X_i-(\theta + \lambda \vec v))^2, \\
			\text{ for }\lambda \in [-\epsilon,0]:  \qquad f(\theta + \lambda \vec v) = f_2(\lambda) &\coloneqq \textstyle\sum_{i \in [n]} \vec b_i^{\T}(X_i-(\theta + \lambda \vec v))^2. 
		\end{align*}
		As $\vec v$ is directed into $\FMset$, $f(\theta+\lambda \vec v)$ is constant for small $\lambda$. Hence $f_1$ must be constant for all $\lambda$. Meanwhile, $-\vec v$ is directed away from $\FMset$ so $f_2$ must be strictly decreasing for $\lambda \in [-\epsilon,0]$. Moreover,
		\begin{align*}
			f_2'(0)  = -\partial_{-\vec v}f(\theta) = 0.
		\end{align*}  
		The quadratic $f_2$ is therefore strictly decreasing on $[-\epsilon,0]$ with $f_2'(0) = 0$, so it follows that $f_2$ is strictly increasing for positive $\lambda$. As $f_1$ is constant for positive $\lambda$ and $f_1(0) = f_2(0) = f(\theta)$, $f_2$ must strictly dominate $f_1$ for small $\lambda$. This contradicts the maximality of $ f_1(\lambda)$ for $\lambda \in [0,\epsilon]$.
	\end{proof}
	
	\begin{lemma}\label{lemma:U_in_ClV}
		For all $\epsilon > 0$, $(X_1,\dots, X_n,\theta) \in U_{F_1,\dots, F_n,\Theta}$, there is some $(Y_1, \dots, Y_n,\theta) \in V_{\CType, \Theta}$ such that $\|X_i - Y_i\|_B < \epsilon$.
	\end{lemma}
	
	\begin{proof}
		We begin by showing that we may assume that $\vec 0 \in \relint \left( \conv \left(\textstyle \bigcup_{i \in [n]} G_i \right) \right)$. Fix some $(\sample, \theta)$ in $U_{\CType,\Theta}$, and assume the contrary, that $\vec 0 \notin \relint \left( \conv \left(\textstyle \bigcup_{i \in [n]} G_i \right) \right)$.
		
		Then there is some $\vec v$ such that for all $i$, $\vec g \in G_i$, we have $\vec g^{\T}\vec v\leq 0$ and there is some $i, \vec g \in G_i$ such that $\vec g^{\T}\vec v< 0$. We can therefore bound the directional derivatives of each term in the Fr\'echet function:
		\begin{align*}
			\partial_{\vec v}(d(\theta,X_i)^2) = \sup_{\vec g \in G_i} [d(\theta, X_i)\vec g^{\T}\vec v] \leq 0, \qquad \text{ and } \qquad
			\partial_{-\vec v}(d(\theta,X_i)^2) = \sup_{\vec g \in G_i} [-d(\theta, X_i)\vec g^{\T}\vec v] \geq 0.
		\end{align*}
		As $\theta$ is a Fr\'echet mean of $X_1, \dots X_n$, every directional derivative of $f$ at $\theta$ is non-negative. That is:
		\begin{align*}
			0 &\leq \partial_{\vec v} f(\theta) = \textstyle \frac{1}{n} \sum_{i \in [n]} \sup_{\vec g \in G_i} [d(\theta, X_i)\vec g^{\T}\vec v] \leq 0.
		\end{align*}
		The $\vec v$-directional derivative of every $d(\theta,X_i)^2$ term is therefore $0$. To compute the directional derivatives explicitly, we note that there is some $\epsilon>0$ and $\vec a_i \in A$ such that, for $\lambda \in [0,\epsilon]$, we have $d(\theta+\lambda \vec v,X_i) = \vec a_i^{\T}(\theta + \lambda \vec v - X_i)$. Then:
		\begin{align*}
			0 &= \partial_{\vec v} (d^2(\theta,X_i)), \\
			&= \left(\frac{\partial}{\partial \lambda}\right)_{\lambda = 0} [\vec a_i^{\T}(\theta + \lambda \vec v - X_i)]^2, \\
			&= (\vec a_i^{\T}\vec v) \cdot d(\theta, X_i).
		\end{align*}
		As discussed in \cref{rmk:positive_dist_to_data}, the fact that the face type is well defined for $\sample$ implies that $d(\theta, X_i)>0$. Therefore we have that $\vec a_i^{\T}\vec v = 0$, and for every $i$:
		\begin{align*}
			d(\theta+\lambda \vec v,X_i) &= \vec a_i^{\T}(\theta + \lambda \vec v - X_i), \\
			&= \vec a_i^{\T}(\theta - X_i), \\
			&= d(\theta,X_i),
		\end{align*}
		which shows that $\theta + \lambda \vec v \in \FMset$ for all $\lambda \in [0,\epsilon]$. 
		By the definition of $U_{\CType,\Theta}$ we know $\theta \in \relint \FMset$, and so $\theta - \lambda \vec v$ must also be in $\FMset$ for sufficiently small $\lambda$. Hence
		\begin{align*}
			0 &= \partial_{-\vec v} f(\theta) =\textstyle \frac{1}{n} \sum_{i \in [n]} \sup_{\vec g \in G_i} [-d(\theta, X_i)\vec g^{\T}\vec v].
		\end{align*}
		The direction $\vec v$ was chosen such that for all $i$, $\vec g \in G_i$, $\vec g^{\T}\vec v\leq 0$, with strict inequality in at least one case. But from the equality above, we conclude that for all $i$, and $\vec g \in G_i$, $-d(\theta, X_i)\vec g^{\T}\vec v = 0$, which contradicts the definition of $\vec v$. Hence we have $\vec 0 \in \relint \left( \conv \left(\textstyle \bigcup_{i \in [n]} G_i \right) \right)$.
		
		By \cite[Theorem 6.9]{rockafellar1970convex}, $\vec 0 \in \relint \left( \conv \left(\textstyle \bigcup_{i \in [n]} G_i \right) \right)$ implies that there exist $\vec g_i' \in \relint G_i$ and $\lambda_i > 0$ such that:
		\[
		\vec 0 = \sum_i \lambda_i \vec g_i'.
		\]
		Also, as $(X_1,\dots, X_n,\theta) \in U_{F_1,\dots, F_n,\Theta}$, we know that $\theta$ is a Fr\'echet mean and $\vec 0 \in \sum_{i \in [n]} \| X_i - \theta \|_B G_i$. Then there also exist $\vec g_i \in G_i, w_i = \|X_i - \theta\|_B > 0$ such that $\vec 0 = \sum_i w_i \vec g_i$.
		Interpolating between these expressions for $\vec 0$ we get:
		\[
		\vec 0 = \sum_i (1-\delta)w_i \vec g_i + \delta \lambda_i \vec g_i' = \sum_i \left((1-\delta)w_i + \delta \lambda_i\right) \left( \frac{(1-\delta)w_i}{(1-\delta)w_i + \delta \lambda_i} \vec g_i + \frac{\delta \lambda_i}{(1-\delta)w_i + \delta \lambda_i} \vec g_i' \right) .
		\]
		As $\vec g_i'$ is in the relative interior of $G_i$, any point on the line segment between $\vec g_i'$ and $\vec g_i$ is also in the relative interior. In particular, for $\delta \in (0,1]$:
		\[
		\frac{(1-\delta)w_i}{(1-\delta)w_i + \delta \lambda_i} \vec g_i + \frac{\delta \lambda_i}{(1-\delta)w_i + \delta \lambda_i} \vec g_i' \eqqcolon \vec g_i'' \in \relint G_i.
		\]
		By picking $\delta >0$ sufficiently small, we can therefore write $\vec 0 = \sum_{i \in [n]} w_i' \vec g_i''$ where $\vec g_i'' \in \relint G_i$ and $w_i' \coloneqq (1-\delta)w_i + \delta \lambda_i$ satisfies $|w_i' - w_i| < \epsilon, w_i'>0$. \\
		We now define $Y_i = \theta + \frac{w_i'}{w_i}(X_i-\theta)$. We will show that $\| Y_i - X_i \|_B < \epsilon$ and $(Y_1, \dots, Y_n,\theta) \in V_{\CType, \Theta}$.
		By the definitions of $Y_i$, $w_i$, and $w_i'$:
		\begin{align*}
			\|Y_i - X_i\|_B &= |w_i'/w_i - 1| \|X_i - \theta\|_B, \\
			&= |w_i' - w_i|, \\
			&< \epsilon.
		\end{align*}
		We also have $\|Y_i - \theta\|_B = w_i'$ and $Y_i - \theta = \tfrac{w_i'}{w_i}(X_i-\theta) \in \relint R_+F_i$ because $(\sample,\theta) \in U_{\CType,\Theta}$.
		By the construction of $w_i', \vec g_i''$:
		\begin{align*}
			0 &\in \textstyle\sum_{i \in [n]} w_i' \relint G_i = \relint \sum_{i \in [n]} w_i' G_i = \relint \sum_{i \in [n]} \|Y_i - \theta\| G_i.
		\end{align*}
		Hence $Y_1,\dots,Y_n,\theta$ is in $V_{\CType,\Theta}$ with $\|X_i-Y_i\|_B < \epsilon$.
	\end{proof}
	
	We use these results to prove \cref{prop:lifted_type_event}.
	
	\begin{proof}[Proof of \cref{prop:lifted_type_event}]
		For each $i$, pick some $\vec a_i$ in $\A(F_i)$. Then for all $X_1,\dots X_k, \theta$ such that $X_i - \theta \in \relint C_i$, we have that $\| X_i -\theta \| = \vec a_i^{\T}(X_i - \theta)$. By the definition of $W$, we conclude that
		\[
		\vec 0 \in \relint \left( \textstyle \sum_{i \in [n]} \|X_i- \theta\|G_i \right) \;\;\text{ if and only if }\;\; \left(\vec a_1^{\T} (X_1-\theta), \dots, \vec a_n^{\T} (X_n-\theta)\right) \in W.
		\]
		Let $L_{\CType}:\R^{k(n+1)}\rightarrow \R^n$ be the linear map $(\sample,\theta) \mapsto \left(\vec a_1^{\T} (X_1-\theta), \dots, \vec a_n^{\T} (X_n-\theta)\right)$, so:
		\[
		V_{\CType, \Theta} =   L_{\CType}^{-1}(W) \; \textstyle\bigcap \; \{ (\sample, \theta) : X_i - \theta \in \relint C_i, \; i = 1,\ldots,n \}.
		\]
		The set $\{(X_1,\ldots,X_n,\theta): X_i - \theta \in \relint C_i \}$ is a relatively open polyhedral cone, as is $L_{\CType}^{-1}(W)$ because $W$ is a relatively open polyhedral cone. So $V_{\CType, \Theta}$ is a relatively open polyhedral cone. 
		
		By Lemma \ref{lemma:RelintGrad_to_RelintFMset}, we deduce that $V_{\CType,\Theta} \subseteq U_{\CType,\Theta}$. Taking $\relint$ and $\closure$ on both sides:
		\begin{align*}
			\closure V_{\CType\Theta} \subseteq \closure U_{\CType\Theta}, \qquad 
			\relint V_{\CType,\Theta} \subseteq \relint U_{\CType,\Theta}.
		\end{align*}
		Lemma \ref{lemma:U_in_ClV} tells us that $U_{\CType,\Theta} \subseteq \closure V_{\CType,\Theta}$. Again, taking $\relint$ and $\closure$ on both sides:
		\begin{align*}
			\closure U_{\CType\Theta} \subseteq \closure V_{\CType\Theta}, \qquad
			\relint U_{\CType,\Theta} \subseteq \relint V_{\CType,\Theta}.
		\end{align*}
		Where the final line is due to the convexity of $V_{\CType,\Theta}$. The result follows.
	\end{proof}
	
\end{document}
